%% Generated by Sphinx.
\def\sphinxdocclass{report}
\documentclass[letterpaper,10pt,english]{sphinxmanual}
\ifdefined\pdfpxdimen
   \let\sphinxpxdimen\pdfpxdimen\else\newdimen\sphinxpxdimen
\fi \sphinxpxdimen=.75bp\relax
\ifdefined\pdfimageresolution
    \pdfimageresolution= \numexpr \dimexpr1in\relax/\sphinxpxdimen\relax
\fi
\newdimen\sphinxremdimen\sphinxremdimen = 10pt
%% let collapsible pdf bookmarks panel have high depth per default
\PassOptionsToPackage{bookmarksdepth=5}{hyperref}

\PassOptionsToPackage{booktabs}{sphinx}
\PassOptionsToPackage{colorrows}{sphinx}

\PassOptionsToPackage{warn}{textcomp}
\usepackage[utf8]{inputenc}
\ifdefined\DeclareUnicodeCharacter
% support both utf8 and utf8x syntaxes
  \ifdefined\DeclareUnicodeCharacterAsOptional
    \def\sphinxDUC#1{\DeclareUnicodeCharacter{"#1}}
  \else
    \let\sphinxDUC\DeclareUnicodeCharacter
  \fi
  \sphinxDUC{00A0}{\nobreakspace}
  \sphinxDUC{2500}{\sphinxunichar{2500}}
  \sphinxDUC{2502}{\sphinxunichar{2502}}
  \sphinxDUC{2514}{\sphinxunichar{2514}}
  \sphinxDUC{251C}{\sphinxunichar{251C}}
  \sphinxDUC{2572}{\textbackslash}
\fi
\usepackage{cmap}
\usepackage[T1]{fontenc}
\usepackage{amsmath,amssymb,amstext}
\usepackage{babel}



\usepackage{tgtermes}
\usepackage{tgheros}
\renewcommand{\ttdefault}{txtt}



\usepackage[Bjarne]{fncychap}
\usepackage{sphinx}

\fvset{fontsize=auto}
\usepackage{geometry}


% Include hyperref last.
\usepackage{hyperref}
% Fix anchor placement for figures with captions.
\usepackage{hypcap}% it must be loaded after hyperref.
% Set up styles of URL: it should be placed after hyperref.
\urlstyle{same}


\usepackage{sphinxmessages}
\setcounter{tocdepth}{1}



\title{Digital Twins Course EDA}
\date{Apr 28, 2026}
\release{0.1}
\author{UTD 7V88 Students}
\newcommand{\sphinxlogo}{\vbox{}}
\renewcommand{\releasename}{Release}
\makeindex
\begin{document}

\ifdefined\shorthandoff
  \ifnum\catcode`\=\string=\active\shorthandoff{=}\fi
  \ifnum\catcode`\"=\active\shorthandoff{"}\fi
\fi

\pagestyle{empty}
\sphinxmaketitle
\pagestyle{plain}
\sphinxtableofcontents
\pagestyle{normal}
\phantomsection\label{\detokenize{index::doc}}


\sphinxAtStartPar
Contents:

\sphinxstepscope


\chapter{Main Execution Script}
\label{\detokenize{simulation_core:main-execution-script}}\label{\detokenize{simulation_core::doc}}
\sphinxAtStartPar
This file serves as the primary entry point for the EDA simulator. It orchestrates
the pipeline: parsing the netlist, constructing the OOP circuit, executing the
numerical engines, and visualizing the results.


\section{Functions}
\label{\detokenize{simulation_core:functions}}\index{built\sphinxhyphen{}in function@\spxentry{built\sphinxhyphen{}in function}!run\_simulation\_core()@\spxentry{run\_simulation\_core()}}\index{run\_simulation\_core()@\spxentry{run\_simulation\_core()}!built\sphinxhyphen{}in function@\spxentry{built\sphinxhyphen{}in function}}

\begin{fulllineitems}
\phantomsection\label{\detokenize{simulation_core:run_simulation_core}}
\pysigstartsignatures
\pysiglinewithargsret
{\sphinxbfcode{\sphinxupquote{run\_simulation\_core}}}
{\sphinxparam{\DUrole{n}{netlist\_path}}\sphinxparamcomma \sphinxparam{\DUrole{n}{output\_nodes}\DUrole{o}{=}\DUrole{default_value}{None}}\sphinxparamcomma \sphinxparam{\DUrole{n}{sensitivity}\DUrole{o}{=}\DUrole{default_value}{False}}\sphinxparamcomma \sphinxparam{\DUrole{n}{global\_adjoint}\DUrole{o}{=}\DUrole{default_value}{False}}\sphinxparamcomma \sphinxparam{\DUrole{n}{keep\_lus}\DUrole{o}{=}\DUrole{default_value}{False}}}
{}
\pysigstopsignatures
\sphinxAtStartPar
Runs the complete SPICE simulation pipeline.
\begin{quote}\begin{description}
\sphinxlineitem{Parameters}\begin{itemize}
\item {} 
\sphinxAtStartPar
\sphinxstyleliteralstrong{\sphinxupquote{netlist\_path}} (\sphinxstyleliteralemphasis{\sphinxupquote{str}}) \textendash{} The file path to the SPICE netlist text file.

\item {} 
\sphinxAtStartPar
\sphinxstyleliteralstrong{\sphinxupquote{output\_nodes}} (\sphinxstyleliteralemphasis{\sphinxupquote{list}}\sphinxstyleliteralemphasis{\sphinxupquote{, }}\sphinxstyleliteralemphasis{\sphinxupquote{optional}}) \textendash{} Target nodes for Adjoint sensitivities. Defaults to None.

\item {} 
\sphinxAtStartPar
\sphinxstyleliteralstrong{\sphinxupquote{sensitivity}} (\sphinxstyleliteralemphasis{\sphinxupquote{bool}}\sphinxstyleliteralemphasis{\sphinxupquote{, }}\sphinxstyleliteralemphasis{\sphinxupquote{optional}}) \textendash{} Toggles the unified continuous sensitivity tensor. Defaults to False.

\item {} 
\sphinxAtStartPar
\sphinxstyleliteralstrong{\sphinxupquote{global\_adjoint}} (\sphinxstyleliteralemphasis{\sphinxupquote{bool}}\sphinxstyleliteralemphasis{\sphinxupquote{, }}\sphinxstyleliteralemphasis{\sphinxupquote{optional}}) \textendash{} Toggles the backward global integral for TRAN. Defaults to False.

\item {} 
\sphinxAtStartPar
\sphinxstyleliteralstrong{\sphinxupquote{keep\_lus}} (\sphinxstyleliteralemphasis{\sphinxupquote{bool}}\sphinxstyleliteralemphasis{\sphinxupquote{, }}\sphinxstyleliteralemphasis{\sphinxupquote{optional}}) \textendash{} Forces LU factorization caching. Defaults to False.

\end{itemize}

\sphinxlineitem{Returns}
\sphinxAtStartPar
A tuple containing \sphinxcode{\sphinxupquote{(Circuit, SimulationResult)}} consisting of the initialized circuit object and the secure data vault with all simulation outputs.

\sphinxlineitem{Return type}
\sphinxAtStartPar
tuple

\end{description}\end{quote}

\end{fulllineitems}


\sphinxstepscope


\chapter{Circuit Representation Module}
\label{\detokenize{circuit:circuit-representation-module}}\label{\detokenize{circuit:circuit-module}}\label{\detokenize{circuit::doc}}
\sphinxAtStartPar
This module provides the central \sphinxcode{\sphinxupquote{Circuit}} class, which acts as the primary data
structure for the simulator. It includes the factory logic to instantiate polymorphic components from parsed netlist dictionaries and generates the Modified Nodal Analysis (MNA) matrix indexing scheme.


\bigskip\hrule\bigskip



\section{Factory Helper}
\label{\detokenize{circuit:factory-helper}}\index{built\sphinxhyphen{}in function@\spxentry{built\sphinxhyphen{}in function}!create\_component()@\spxentry{create\_component()}}\index{create\_component()@\spxentry{create\_component()}!built\sphinxhyphen{}in function@\spxentry{built\sphinxhyphen{}in function}}

\begin{fulllineitems}
\phantomsection\label{\detokenize{circuit:create_component}}
\pysigstartsignatures
\pysiglinewithargsret
{\sphinxbfcode{\sphinxupquote{create\_component}}}
{\sphinxparam{\DUrole{n}{name}}\sphinxparamcomma \sphinxparam{\DUrole{n}{data\_dict}}}
{}
\pysigstopsignatures
\sphinxAtStartPar
Instantiates the correct Component subclass based on the netlist type.
\begin{quote}\begin{description}
\sphinxlineitem{Parameters}\begin{itemize}
\item {} 
\sphinxAtStartPar
\sphinxstyleliteralstrong{\sphinxupquote{name}} (\sphinxstyleliteralemphasis{\sphinxupquote{str}}) \textendash{} The unique netlist name of the component (e.g., \sphinxcode{\sphinxupquote{\textquotesingle{}R1\textquotesingle{}}} or \sphinxcode{\sphinxupquote{\textquotesingle{}M\_MAIN\textquotesingle{}}}).

\item {} 
\sphinxAtStartPar
\sphinxstyleliteralstrong{\sphinxupquote{data\_dict}} (\sphinxstyleliteralemphasis{\sphinxupquote{dict}}) \textendash{} The parsed parameter dictionary containing at minimum a \sphinxcode{\sphinxupquote{\textquotesingle{}type\textquotesingle{}}} key
(e.g. \sphinxcode{\sphinxupquote{\{\textquotesingle{}type\textquotesingle{}: \textquotesingle{}R\textquotesingle{}, \textquotesingle{}n1\textquotesingle{}: \textquotesingle{}1\textquotesingle{}, \textquotesingle{}n2\textquotesingle{}: \textquotesingle{}0\textquotesingle{}, \textquotesingle{}value\textquotesingle{}: 1000\}}}).

\end{itemize}

\sphinxlineitem{Returns}
\sphinxAtStartPar
An instantiated component object.

\sphinxlineitem{Return type}
\sphinxAtStartPar
{\hyperref[\detokenize{component_base:Component}]{\sphinxcrossref{Component}}}

\sphinxlineitem{Raises}\begin{itemize}
\item {} 
\sphinxAtStartPar
\sphinxstyleliteralstrong{\sphinxupquote{KeyError}} \textendash{} If the \sphinxcode{\sphinxupquote{\textquotesingle{}type\textquotesingle{}}} key is missing.

\item {} 
\sphinxAtStartPar
\sphinxstyleliteralstrong{\sphinxupquote{ValueError}} \textendash{} If the component type is not recognized.

\end{itemize}

\end{description}\end{quote}

\end{fulllineitems}



\bigskip\hrule\bigskip



\section{Circuit Class}
\label{\detokenize{circuit:circuit-class}}\index{Circuit (built\sphinxhyphen{}in class)@\spxentry{Circuit}\spxextra{built\sphinxhyphen{}in class}}

\begin{fulllineitems}
\phantomsection\label{\detokenize{circuit:Circuit}}
\pysigstartsignatures
\pysiglinewithargsret
{\sphinxbfcode{\sphinxupquote{\DUrole{k}{class}\DUrole{w}{ }}}\sphinxbfcode{\sphinxupquote{Circuit}}}
{\sphinxparam{\DUrole{n}{parsed\_components}}}
{}
\pysigstopsignatures
\sphinxAtStartPar
Represents the physical circuit, containing components and matrix topology.
\index{components (Circuit attribute)@\spxentry{components}\spxextra{Circuit attribute}}

\begin{fulllineitems}
\phantomsection\label{\detokenize{circuit:Circuit.components}}
\pysigstartsignatures
\pysigline
{\sphinxbfcode{\sphinxupquote{components}}}
\pysigstopsignatures
\sphinxAtStartPar
A list of instantiated polymorphic component objects.

\end{fulllineitems}

\index{components\_dict (Circuit attribute)@\spxentry{components\_dict}\spxextra{Circuit attribute}}

\begin{fulllineitems}
\phantomsection\label{\detokenize{circuit:Circuit.components_dict}}
\pysigstartsignatures
\pysigline
{\sphinxbfcode{\sphinxupquote{components\_dict}}}
\pysigstopsignatures
\sphinxAtStartPar
A name‑to‑object lookup dictionary for fast access.

\end{fulllineitems}

\index{node\_map (Circuit attribute)@\spxentry{node\_map}\spxextra{Circuit attribute}}

\begin{fulllineitems}
\phantomsection\label{\detokenize{circuit:Circuit.node_map}}
\pysigstartsignatures
\pysigline
{\sphinxbfcode{\sphinxupquote{node\_map}}}
\pysigstopsignatures
\sphinxAtStartPar
Maps node names (and MNA branch names) to integer matrix indices.

\end{fulllineitems}

\index{total\_dim (Circuit attribute)@\spxentry{total\_dim}\spxextra{Circuit attribute}}

\begin{fulllineitems}
\phantomsection\label{\detokenize{circuit:Circuit.total_dim}}
\pysigstartsignatures
\pysigline
{\sphinxbfcode{\sphinxupquote{total\_dim}}}
\pysigstopsignatures
\sphinxAtStartPar
The total dimension of the MNA matrix (N × N).

\end{fulllineitems}

\index{get\_idx() (Circuit method)@\spxentry{get\_idx()}\spxextra{Circuit method}}

\begin{fulllineitems}
\phantomsection\label{\detokenize{circuit:Circuit.get_idx}}
\pysigstartsignatures
\pysiglinewithargsret
{\sphinxbfcode{\sphinxupquote{get\_idx}}}
{\sphinxparam{\DUrole{n}{node}}}
{}
\pysigstopsignatures
\sphinxAtStartPar
Retrieves the matrix row/column index for a given node.
\begin{quote}\begin{description}
\sphinxlineitem{Parameters}
\sphinxAtStartPar
\sphinxstyleliteralstrong{\sphinxupquote{node}} \textendash{} Ground nodes (\sphinxcode{\sphinxupquote{0}}, \sphinxcode{\sphinxupquote{"0"}}, or \sphinxcode{\sphinxupquote{"GND"}}) return \sphinxcode{\sphinxupquote{None}}.

\sphinxlineitem{Returns}
\sphinxAtStartPar
Integer index or \sphinxcode{\sphinxupquote{None}}.

\end{description}\end{quote}

\end{fulllineitems}

\index{\_build\_node\_index() (Circuit method)@\spxentry{\_build\_node\_index()}\spxextra{Circuit method}}

\begin{fulllineitems}
\phantomsection\label{\detokenize{circuit:Circuit._build_node_index}}
\pysigstartsignatures
\pysiglinewithargsret
{\sphinxbfcode{\sphinxupquote{\_build\_node\_index}}}
{}
{}
\pysigstopsignatures
\sphinxAtStartPar
Scans components to build the MNA matrix coordinate map.

\end{fulllineitems}

\index{get\_component() (Circuit method)@\spxentry{get\_component()}\spxextra{Circuit method}}

\begin{fulllineitems}
\phantomsection\label{\detokenize{circuit:Circuit.get_component}}
\pysigstartsignatures
\pysiglinewithargsret
{\sphinxbfcode{\sphinxupquote{get\_component}}}
{\sphinxparam{\DUrole{n}{name}}}
{}
\pysigstopsignatures
\sphinxAtStartPar
Retrieves a component object by its netlist name.
\begin{quote}\begin{description}
\sphinxlineitem{Parameters}
\sphinxAtStartPar
\sphinxstyleliteralstrong{\sphinxupquote{name}} (\sphinxstyleliteralemphasis{\sphinxupquote{str}}) \textendash{} The component’s netlist identifier.

\sphinxlineitem{Raises}
\sphinxAtStartPar
\sphinxstyleliteralstrong{\sphinxupquote{KeyError}} \textendash{} If the component is not found.

\end{description}\end{quote}

\end{fulllineitems}

\index{differentiable\_params (Circuit attribute)@\spxentry{differentiable\_params}\spxextra{Circuit attribute}}

\begin{fulllineitems}
\phantomsection\label{\detokenize{circuit:Circuit.differentiable_params}}
\pysigstartsignatures
\pysigline
{\sphinxbfcode{\sphinxupquote{differentiable\_params}}}
\pysigstopsignatures
\sphinxAtStartPar
A master list of all tunable parameters in the circuit.  It aggregates
\sphinxcode{\sphinxupquote{differentiable\_params}} from each individual
component, ensuring no duplicates.

\end{fulllineitems}


\end{fulllineitems}


\sphinxstepscope


\chapter{Thermal Voltage Helper}
\label{\detokenize{thermal_constants:thermal-voltage-helper}}\label{\detokenize{thermal_constants:thermal-constants}}\label{\detokenize{thermal_constants::doc}}
\sphinxAtStartPar
This tiny module simply pulls constants from \sphinxcode{\sphinxupquote{scipy.constants}} and
computes the thermal voltage at room temperature (298 K).  It is handy for
any SPICE‑style circuit simulation that needs \sphinxcode{\sphinxupquote{Vt}}.
\index{e (built\sphinxhyphen{}in variable)@\spxentry{e}\spxextra{built\sphinxhyphen{}in variable}}

\begin{fulllineitems}
\phantomsection\label{\detokenize{thermal_constants:e}}
\pysigstartsignatures
\pysigline
{\sphinxbfcode{\sphinxupquote{e}}\sphinxbfcode{\sphinxupquote{\DUrole{w}{ }float}}}
\pysigstopsignatures
\sphinxAtStartPar
Elementary charge in coulombs \textendash{} imported as \sphinxcode{\sphinxupquote{c.e}} from SciPy.

\end{fulllineitems}

\index{kb (built\sphinxhyphen{}in variable)@\spxentry{kb}\spxextra{built\sphinxhyphen{}in variable}}

\begin{fulllineitems}
\phantomsection\label{\detokenize{thermal_constants:kb}}
\pysigstartsignatures
\pysigline
{\sphinxbfcode{\sphinxupquote{kb}}\sphinxbfcode{\sphinxupquote{\DUrole{w}{ }float}}}
\pysigstopsignatures
\sphinxAtStartPar
Boltzmann constant (J K⁻\(\sp{\text{1}}\)) \textendash{} imported as \sphinxcode{\sphinxupquote{c.k}}.

\end{fulllineitems}

\index{T (built\sphinxhyphen{}in variable)@\spxentry{T}\spxextra{built\sphinxhyphen{}in variable}}

\begin{fulllineitems}
\phantomsection\label{\detokenize{thermal_constants:T}}
\pysigstartsignatures
\pysigline
{\sphinxbfcode{\sphinxupquote{T}}\sphinxbfcode{\sphinxupquote{\DUrole{w}{ }float}}}
\pysigstopsignatures
\sphinxAtStartPar
Reference temperature in kelvin (default 298.15 K).

\end{fulllineitems}

\index{Vt (built\sphinxhyphen{}in variable)@\spxentry{Vt}\spxextra{built\sphinxhyphen{}in variable}}

\begin{fulllineitems}
\phantomsection\label{\detokenize{thermal_constants:Vt}}
\pysigstartsignatures
\pysigline
{\sphinxbfcode{\sphinxupquote{Vt}}\sphinxbfcode{\sphinxupquote{\DUrole{w}{ }float}}}
\pysigstopsignatures
\sphinxAtStartPar
Thermal voltage, computed as \sphinxcode{\sphinxupquote{(kb * T) / e}} (\textasciitilde{}26 mV at room temperature).

\end{fulllineitems}



\section{Usage example}
\label{\detokenize{thermal_constants:usage-example}}
\begin{sphinxVerbatim}[commandchars=\\\{\}]
\PYG{g+gp}{\PYGZgt{}\PYGZgt{}\PYGZgt{} }\PYG{k+kn}{from}\PYG{+w}{ }\PYG{n+nn}{thermal\PYGZus{}constants}\PYG{+w}{ }\PYG{k+kn}{import} \PYG{n}{Vt}
\PYG{g+gp}{\PYGZgt{}\PYGZgt{}\PYGZgt{} }\PYG{n}{Vt}
\PYG{g+go}{0.025851999...}
\end{sphinxVerbatim}

\sphinxAtStartPar
The module is intentionally lightweight; if you need a different temperature,
edit {\hyperref[\detokenize{thermal_constants:T}]{\sphinxcrossref{\sphinxcode{\sphinxupquote{T}}}}} or recompute {\hyperref[\detokenize{thermal_constants:Vt}]{\sphinxcrossref{\sphinxcode{\sphinxupquote{Vt}}}}} manually.

\sphinxstepscope


\chapter{Diode \& NMOS Level‑1 Models}
\label{\detokenize{diode_mosfet_models:diode-nmos-level1-models}}\label{\detokenize{diode_mosfet_models:diode-and-nmos-models}}\label{\detokenize{diode_mosfet_models::doc}}
\sphinxAtStartPar
This module implements two small‑signal models that are useful in SPICE‑style circuit simulation:
\begin{itemize}
\item {} 
\sphinxAtStartPar
{\hyperref[\detokenize{diode_mosfet_models:evaluate_diode}]{\sphinxcrossref{\sphinxcode{\sphinxupquote{evaluate\_diode()}}}}} \textendash{} Shockley diode model with sensitivity to the saturation current.

\item {} 
\sphinxAtStartPar
{\hyperref[\detokenize{diode_mosfet_models:evaluate_nmos}]{\sphinxcrossref{\sphinxcode{\sphinxupquote{evaluate\_nmos()}}}}} \textendash{} Level 1 MOSFET (NMOS) model including gradients for \sphinxstyleemphasis{V\textless{}sub\textgreater{}TO\textless{}/sub\textgreater{}} and \sphinxstyleemphasis{β\textless{}sub\textgreater{}n\textless{}/sub\textgreater{}}.

\end{itemize}

\sphinxAtStartPar
The functions return dictionaries containing the large‑signal quantities as well as the required small‑signal conductances and sensitivity gradients.


\bigskip\hrule\bigskip

\index{built\sphinxhyphen{}in function@\spxentry{built\sphinxhyphen{}in function}!evaluate\_diode()@\spxentry{evaluate\_diode()}}\index{evaluate\_diode()@\spxentry{evaluate\_diode()}!built\sphinxhyphen{}in function@\spxentry{built\sphinxhyphen{}in function}}

\begin{fulllineitems}
\phantomsection\label{\detokenize{diode_mosfet_models:evaluate_diode}}
\pysigstartsignatures
\pysiglinewithargsret
{\sphinxbfcode{\sphinxupquote{evaluate\_diode}}}
{\sphinxparam{\DUrole{n}{vd}}\sphinxparamcomma \sphinxparam{\DUrole{n}{Is}}\sphinxparamcomma \sphinxparam{\DUrole{n}{Vt}}}
{}
\pysigstopsignatures
\sphinxAtStartPar
Evaluates the Shockley diode model.
\begin{quote}\begin{description}
\sphinxlineitem{Parameters}\begin{itemize}
\item {} 
\sphinxAtStartPar
\sphinxstyleliteralstrong{\sphinxupquote{vd}} (\sphinxstyleliteralemphasis{\sphinxupquote{float}}) \textendash{} Voltage across the diode (Anode \textendash{} Cathode).

\item {} 
\sphinxAtStartPar
\sphinxstyleliteralstrong{\sphinxupquote{Is}} (\sphinxstyleliteralemphasis{\sphinxupquote{float}}) \textendash{} Saturation current (e.g., 1e‑14 A).

\item {} 
\sphinxAtStartPar
\sphinxstyleliteralstrong{\sphinxupquote{Vt}} (\sphinxstyleliteralemphasis{\sphinxupquote{float}}) \textendash{} Thermal voltage (e.g., 0.02585 V).

\end{itemize}

\sphinxlineitem{Returns}
\sphinxAtStartPar

\sphinxAtStartPar
\sphinxcode{\sphinxupquote{dict}} with keys:
\begin{itemize}
\item {} 
\sphinxAtStartPar
\sphinxcode{\sphinxupquote{"I\_D"}} \textendash{} Large‑signal current in amperes.

\item {} 
\sphinxAtStartPar
\sphinxcode{\sphinxupquote{"gd"}}   \textendash{} Small‑signal conductance dI/dv (siemens).

\item {} 
\sphinxAtStartPar
\sphinxcode{\sphinxupquote{"dId\_dIs"}}
\textendash{} Sensitivity gradient ∂I/∂Is.

\end{itemize}


\sphinxlineitem{Return type}
\sphinxAtStartPar
dict

\end{description}\end{quote}

\end{fulllineitems}



\bigskip\hrule\bigskip

\index{built\sphinxhyphen{}in function@\spxentry{built\sphinxhyphen{}in function}!evaluate\_nmos()@\spxentry{evaluate\_nmos()}}\index{evaluate\_nmos()@\spxentry{evaluate\_nmos()}!built\sphinxhyphen{}in function@\spxentry{built\sphinxhyphen{}in function}}

\begin{fulllineitems}
\phantomsection\label{\detokenize{diode_mosfet_models:evaluate_nmos}}
\pysigstartsignatures
\pysiglinewithargsret
{\sphinxbfcode{\sphinxupquote{evaluate\_nmos}}}
{\sphinxparam{\DUrole{n}{vgs}}\sphinxparamcomma \sphinxparam{\DUrole{n}{vds}}\sphinxparamcomma \sphinxparam{\DUrole{n}{VTO}}\sphinxparamcomma \sphinxparam{\DUrole{n}{Bn}}}
{}
\pysigstopsignatures
\sphinxAtStartPar
Evaluates the Level 1 NMOS model with sensitivity gradients.
\begin{quote}\begin{description}
\sphinxlineitem{Parameters}\begin{itemize}
\item {} 
\sphinxAtStartPar
\sphinxstyleliteralstrong{\sphinxupquote{vgs}} (\sphinxstyleliteralemphasis{\sphinxupquote{float}}) \textendash{} Gate‑source voltage.

\item {} 
\sphinxAtStartPar
\sphinxstyleliteralstrong{\sphinxupquote{vds}} (\sphinxstyleliteralemphasis{\sphinxupquote{float}}) \textendash{} Drain‑source voltage.

\item {} 
\sphinxAtStartPar
\sphinxstyleliteralstrong{\sphinxupquote{VTO}} (\sphinxstyleliteralemphasis{\sphinxupquote{float}}) \textendash{} Threshold voltage (V\textless{}sub\textgreater{}TO\textless{}/sub\textgreater{}).

\item {} 
\sphinxAtStartPar
\sphinxstyleliteralstrong{\sphinxupquote{Bn}} (\sphinxstyleliteralemphasis{\sphinxupquote{float}}) \textendash{} Transconductance parameter β\textless{}sub\textgreater{}n\textless{}/sub\textgreater{} = μₙC\textless{}sub\textgreater{}ox\textless{}/sub\textgreater{}(W/L).

\end{itemize}

\sphinxlineitem{Returns}
\sphinxAtStartPar

\sphinxAtStartPar
\sphinxcode{\sphinxupquote{dict}} with keys:
\begin{itemize}
\item {} 
\sphinxAtStartPar
\sphinxcode{\sphinxupquote{"I\_D"}}   \textendash{} Drain current (A).

\item {} 
\sphinxAtStartPar
\sphinxcode{\sphinxupquote{"gm"}}    \textendash{} Transconductance dI/dV\textless{}sub\textgreater{}GS\textless{}/sub\textgreater{}.

\item {} 
\sphinxAtStartPar
\sphinxcode{\sphinxupquote{"gds"}}   \textendash{} Output conductance dI/dV\textless{}sub\textgreater{}DS\textless{}/sub\textgreater{}.

\item {} 
\sphinxAtStartPar
\sphinxcode{\sphinxupquote{"dId\_dBn"}}
\textendash{} Sensitivity ∂I/∂β\textless{}sub\textgreater{}n\textless{}/sub\textgreater{}.

\item {} 
\sphinxAtStartPar
\sphinxcode{\sphinxupquote{"dId\_dVTO"}}
\textendash{} Sensitivity ∂I/∂V\textless{}sub\textgreater{}TO\textless{}/sub\textgreater{}.

\end{itemize}


\sphinxlineitem{Return type}
\sphinxAtStartPar
dict

\end{description}\end{quote}

\end{fulllineitems}


\sphinxstepscope


\chapter{Simulation Result \& Sensitivity Data}
\label{\detokenize{results:simulation-result-sensitivity-data}}\label{\detokenize{results:results-module}}\label{\detokenize{results::doc}}
\sphinxAtStartPar
This module provides two high‑level containers that are returned by the simulator:
{\color{red}\bfseries{}\textasciigrave{}\textasciigrave{}}{\color{red}\bfseries{}\textasciigrave{}}
\begin{itemize}
\item {} 
\sphinxAtStartPar
{\hyperref[\detokenize{results:SensitivityData}]{\sphinxcrossref{\sphinxcode{\sphinxupquote{SensitivityData}}}}} \textendash{} A unified 3‑D tensor for all sensitivity results.

\item {} 
\sphinxAtStartPar
{\hyperref[\detokenize{results:SimulationResult}]{\sphinxcrossref{\sphinxcode{\sphinxupquote{SimulationResult}}}}} \textendash{} The “vault” that stores raw simulation outputs and, when available, the associated sensitivities.

\end{itemize}


\bigskip\hrule\bigskip

\index{SensitivityData (built\sphinxhyphen{}in class)@\spxentry{SensitivityData}\spxextra{built\sphinxhyphen{}in class}}

\begin{fulllineitems}
\phantomsection\label{\detokenize{results:SensitivityData}}
\pysigstartsignatures
\pysiglinewithargsret
{\sphinxbfcode{\sphinxupquote{\DUrole{k}{class}\DUrole{w}{ }}}\sphinxbfcode{\sphinxupquote{SensitivityData}}}
{\sphinxparam{\DUrole{n}{sweep\_axis}}\sphinxparamcomma \sphinxparam{\DUrole{n}{param\_names}}\sphinxparamcomma \sphinxparam{\DUrole{n}{output\_nodes}}\sphinxparamcomma \sphinxparam{\DUrole{n}{domain}\DUrole{o}{=}\DUrole{default_value}{\textquotesingle{}static\textquotesingle{}}}}
{}
\pysigstopsignatures
\sphinxAtStartPar
Holds a 3‑D tensor of shape \sphinxstyleemphasis{(Parameters × Output Nodes × Sweep Steps)}.
The class offers convenient query helpers that hide the raw NumPy indexing from the user.
\index{sweep\_axis (SensitivityData attribute)@\spxentry{sweep\_axis}\spxextra{SensitivityData attribute}}

\begin{fulllineitems}
\phantomsection\label{\detokenize{results:SensitivityData.sweep_axis}}
\pysigstartsignatures
\pysigline
{\sphinxbfcode{\sphinxupquote{sweep\_axis}}}
\pysigstopsignatures
\sphinxAtStartPar
\sphinxcode{\sphinxupquote{numpy.ndarray}} \textendash{} Independent variable array (time, frequency, voltage, or single value for .OP).

\end{fulllineitems}

\index{domain (SensitivityData attribute)@\spxentry{domain}\spxextra{SensitivityData attribute}}

\begin{fulllineitems}
\phantomsection\label{\detokenize{results:SensitivityData.domain}}
\pysigstartsignatures
\pysigline
{\sphinxbfcode{\sphinxupquote{domain}}}
\pysigstopsignatures
\sphinxAtStartPar
Physical domain of the sweep (\sphinxtitleref{“time”}, \sphinxtitleref{“frequency”}, \sphinxtitleref{“voltage”}, \sphinxtitleref{“static”}).

\end{fulllineitems}

\index{param\_names (SensitivityData attribute)@\spxentry{param\_names}\spxextra{SensitivityData attribute}}

\begin{fulllineitems}
\phantomsection\label{\detokenize{results:SensitivityData.param_names}}
\pysigstartsignatures
\pysigline
{\sphinxbfcode{\sphinxupquote{param\_names}}}
\pysigstopsignatures
\sphinxAtStartPar
Ordered list of parameter names (axis‑0 labels).

\end{fulllineitems}

\index{output\_nodes (SensitivityData attribute)@\spxentry{output\_nodes}\spxextra{SensitivityData attribute}}

\begin{fulllineitems}
\phantomsection\label{\detokenize{results:SensitivityData.output_nodes}}
\pysigstartsignatures
\pysigline
{\sphinxbfcode{\sphinxupquote{output\_nodes}}}
\pysigstopsignatures
\sphinxAtStartPar
Ordered list of output node names (axis‑1 labels).

\end{fulllineitems}

\index{data (SensitivityData attribute)@\spxentry{data}\spxextra{SensitivityData attribute}}

\begin{fulllineitems}
\phantomsection\label{\detokenize{results:SensitivityData.data}}
\pysigstartsignatures
\pysigline
{\sphinxbfcode{\sphinxupquote{data}}}
\pysigstopsignatures
\sphinxAtStartPar
\sphinxcode{\sphinxupquote{numpy.ndarray}} \textendash{} The underlying 3‑D sensitivity tensor.

\end{fulllineitems}

\index{param\_index (SensitivityData attribute)@\spxentry{param\_index}\spxextra{SensitivityData attribute}}

\begin{fulllineitems}
\phantomsection\label{\detokenize{results:SensitivityData.param_index}}
\pysigstartsignatures
\pysigline
{\sphinxbfcode{\sphinxupquote{param\_index}}}
\pysigstopsignatures
\end{fulllineitems}

\index{output\_index (SensitivityData attribute)@\spxentry{output\_index}\spxextra{SensitivityData attribute}}

\begin{fulllineitems}
\phantomsection\label{\detokenize{results:SensitivityData.output_index}}
\pysigstartsignatures
\pysigline
{\sphinxbfcode{\sphinxupquote{output\_index}}}
\pysigstopsignatures
\sphinxAtStartPar
Dictionaries mapping names \(\rightarrow\) indices for fast lookup.

\end{fulllineitems}

\index{get\_sweep\_series() (SensitivityData method)@\spxentry{get\_sweep\_series()}\spxextra{SensitivityData method}}

\begin{fulllineitems}
\phantomsection\label{\detokenize{results:SensitivityData.get_sweep_series}}
\pysigstartsignatures
\pysiglinewithargsret
{\sphinxbfcode{\sphinxupquote{get\_sweep\_series}}}
{\sphinxparam{\DUrole{n}{param}}\sphinxparamcomma \sphinxparam{\DUrole{n}{output\_node}}}
{}
\pysigstopsignatures
\sphinxAtStartPar
Returns the 1‑D array of partial derivatives \sphinxstyleemphasis{∂output/∂param} across the sweep axis.
\begin{quote}\begin{description}
\sphinxlineitem{Parameters}\begin{itemize}
\item {} 
\sphinxAtStartPar
\sphinxstyleliteralstrong{\sphinxupquote{param}} (\sphinxstyleliteralemphasis{\sphinxupquote{str}}) \textendash{} Parameter name.

\item {} 
\sphinxAtStartPar
\sphinxstyleliteralstrong{\sphinxupquote{output\_node}} (\sphinxstyleliteralemphasis{\sphinxupquote{str}}\sphinxstyleliteralemphasis{\sphinxupquote{|}}\sphinxstyleliteralemphasis{\sphinxupquote{int}}) \textendash{} Target node.

\end{itemize}

\sphinxlineitem{Returns}
\sphinxAtStartPar
\sphinxcode{\sphinxupquote{numpy.ndarray}} \textendash{} 1‑D sensitivity series.

\end{description}\end{quote}

\end{fulllineitems}

\index{get\_matrix\_at\_step() (SensitivityData method)@\spxentry{get\_matrix\_at\_step()}\spxextra{SensitivityData method}}

\begin{fulllineitems}
\phantomsection\label{\detokenize{results:SensitivityData.get_matrix_at_step}}
\pysigstartsignatures
\pysiglinewithargsret
{\sphinxbfcode{\sphinxupquote{get\_matrix\_at\_step}}}
{\sphinxparam{\DUrole{n}{step\_idx}}}
{}
\pysigstopsignatures
\sphinxAtStartPar
Returns a 2‑D matrix \sphinxstyleemphasis{(n\_params × n\_outputs)} at the given sweep index.
\begin{quote}\begin{description}
\sphinxlineitem{Parameters}
\sphinxAtStartPar
\sphinxstyleliteralstrong{\sphinxupquote{step\_idx}} (\sphinxstyleliteralemphasis{\sphinxupquote{int}}) \textendash{} Sweep index.

\sphinxlineitem{Returns}
\sphinxAtStartPar
\sphinxcode{\sphinxupquote{numpy.ndarray}}.

\end{description}\end{quote}

\end{fulllineitems}

\index{print\_matrix\_at\_step() (SensitivityData method)@\spxentry{print\_matrix\_at\_step()}\spxextra{SensitivityData method}}

\begin{fulllineitems}
\phantomsection\label{\detokenize{results:SensitivityData.print_matrix_at_step}}
\pysigstartsignatures
\pysiglinewithargsret
{\sphinxbfcode{\sphinxupquote{print\_matrix\_at\_step}}}
{\sphinxparam{\DUrole{n}{step\_idx}}\sphinxparamcomma \sphinxparam{\DUrole{n}{output\_node}}\sphinxparamcomma \sphinxparam{\DUrole{n}{sens\_parameter}}}
{}
\pysigstopsignatures
\sphinxAtStartPar
Nicely formatted console output of the sensitivity matrix for a single
sweep step.  Useful for quick debugging or CLI reporting.

\end{fulllineitems}


\end{fulllineitems}



\bigskip\hrule\bigskip

\index{SimulationResult (built\sphinxhyphen{}in class)@\spxentry{SimulationResult}\spxextra{built\sphinxhyphen{}in class}}

\begin{fulllineitems}
\phantomsection\label{\detokenize{results:SimulationResult}}
\pysigstartsignatures
\pysiglinewithargsret
{\sphinxbfcode{\sphinxupquote{\DUrole{k}{class}\DUrole{w}{ }}}\sphinxbfcode{\sphinxupquote{SimulationResult}}}
{\sphinxparam{\DUrole{n}{analysis\_type}}\sphinxparamcomma \sphinxparam{\DUrole{n}{sweep\_axis}}\sphinxparamcomma \sphinxparam{\DUrole{n}{VI\_matrix}}\sphinxparamcomma \sphinxparam{\DUrole{n}{node\_map}}}
{}
\pysigstopsignatures
\sphinxAtStartPar
A container that holds all information produced by a simulation run:
\begin{itemize}
\item {} 
\sphinxAtStartPar
The raw solution matrix (\sphinxtitleref{VI}).

\item {} 
\sphinxAtStartPar
The node/branch mapping.

\item {} 
\sphinxAtStartPar
Optional {\hyperref[\detokenize{results:SimulationResult.sensitivities}]{\sphinxcrossref{\sphinxcode{\sphinxupquote{sensitivities}}}}} (an instance of {\hyperref[\detokenize{results:SensitivityData}]{\sphinxcrossref{\sphinxcode{\sphinxupquote{SensitivityData}}}}}).

\end{itemize}
\index{type (SimulationResult attribute)@\spxentry{type}\spxextra{SimulationResult attribute}}

\begin{fulllineitems}
\phantomsection\label{\detokenize{results:SimulationResult.type}}
\pysigstartsignatures
\pysigline
{\sphinxbfcode{\sphinxupquote{type}}}
\pysigstopsignatures
\sphinxAtStartPar
Analysis type string (\sphinxtitleref{‘.TRAN’}, \sphinxtitleref{‘.AC’}, \sphinxtitleref{‘.DC’}, \sphinxtitleref{‘.OP’}).

\end{fulllineitems}

\index{sweep\_axis (SimulationResult attribute)@\spxentry{sweep\_axis}\spxextra{SimulationResult attribute}}

\begin{fulllineitems}
\phantomsection\label{\detokenize{results:SimulationResult.sweep_axis}}
\pysigstartsignatures
\pysigline
{\sphinxbfcode{\sphinxupquote{sweep\_axis}}}
\pysigstopsignatures
\sphinxAtStartPar
Independent variable array.

\end{fulllineitems}

\index{VI (SimulationResult attribute)@\spxentry{VI}\spxextra{SimulationResult attribute}}

\begin{fulllineitems}
\phantomsection\label{\detokenize{results:SimulationResult.VI}}
\pysigstartsignatures
\pysigline
{\sphinxbfcode{\sphinxupquote{VI}}}
\pysigstopsignatures
\sphinxAtStartPar
Raw solver output matrix (2‑D for sweeps, 1‑D for .OP).

\end{fulllineitems}

\index{node\_map (SimulationResult attribute)@\spxentry{node\_map}\spxextra{SimulationResult attribute}}

\begin{fulllineitems}
\phantomsection\label{\detokenize{results:SimulationResult.node_map}}
\pysigstartsignatures
\pysigline
{\sphinxbfcode{\sphinxupquote{node\_map}}}
\pysigstopsignatures
\sphinxAtStartPar
Mapping of node/branch names \(\rightarrow\) column indices in {\hyperref[\detokenize{results:SimulationResult.VI}]{\sphinxcrossref{\sphinxcode{\sphinxupquote{VI}}}}}.

\end{fulllineitems}

\index{sensitivities (SimulationResult attribute)@\spxentry{sensitivities}\spxextra{SimulationResult attribute}}

\begin{fulllineitems}
\phantomsection\label{\detokenize{results:SimulationResult.sensitivities}}
\pysigstartsignatures
\pysigline
{\sphinxbfcode{\sphinxupquote{sensitivities}}}
\pysigstopsignatures
\sphinxAtStartPar
Optional {\hyperref[\detokenize{results:SensitivityData}]{\sphinxcrossref{\sphinxcode{\sphinxupquote{SensitivityData}}}}} instance.

\end{fulllineitems}

\index{get\_vhats() (SimulationResult method)@\spxentry{get\_vhats()}\spxextra{SimulationResult method}}

\begin{fulllineitems}
\phantomsection\label{\detokenize{results:SimulationResult.get_vhats}}
\pysigstartsignatures
\pysiglinewithargsret
{\sphinxbfcode{\sphinxupquote{get\_vhats}}}
{\sphinxparam{\DUrole{n}{node}}}
{}
\pysigstopsignatures
\sphinxAtStartPar
Retrieves the raw adjoint vector history (useful for fault analysis).

\end{fulllineitems}

\index{\_resolve\_node\_key() (SimulationResult method)@\spxentry{\_resolve\_node\_key()}\spxextra{SimulationResult method}}

\begin{fulllineitems}
\phantomsection\label{\detokenize{results:SimulationResult._resolve_node_key}}
\pysigstartsignatures
\pysiglinewithargsret
{\sphinxbfcode{\sphinxupquote{\_resolve\_node\_key}}}
{\sphinxparam{\DUrole{n}{node}}}
{}
\pysigstopsignatures
\sphinxAtStartPar
Internal helper that normalises user‑supplied node identifiers
to the exact key stored in {\hyperref[\detokenize{results:SimulationResult.node_map}]{\sphinxcrossref{\sphinxcode{\sphinxupquote{node\_map}}}}}.

\end{fulllineitems}

\index{get\_voltage() (SimulationResult method)@\spxentry{get\_voltage()}\spxextra{SimulationResult method}}

\begin{fulllineitems}
\phantomsection\label{\detokenize{results:SimulationResult.get_voltage}}
\pysigstartsignatures
\pysiglinewithargsret
{\sphinxbfcode{\sphinxupquote{get\_voltage}}}
{\sphinxparam{\DUrole{n}{node}}}
{}
\pysigstopsignatures
\sphinxAtStartPar
Returns the voltage or branch current array for \sphinxstyleemphasis{node}.
Handles complex values for .AC, real values otherwise.

\end{fulllineitems}

\index{get\_sensitivity\_parameters() (SimulationResult method)@\spxentry{get\_sensitivity\_parameters()}\spxextra{SimulationResult method}}

\begin{fulllineitems}
\phantomsection\label{\detokenize{results:SimulationResult.get_sensitivity_parameters}}
\pysigstartsignatures
\pysiglinewithargsret
{\sphinxbfcode{\sphinxupquote{get\_sensitivity\_parameters}}}
{\sphinxparam{\DUrole{n}{node}}}
{}
\pysigstopsignatures
\sphinxAtStartPar
List of parameters that have sensitivity data for \sphinxstyleemphasis{node}.
Useful for populating UI dropdowns.

\end{fulllineitems}

\index{get\_sensitivity() (SimulationResult method)@\spxentry{get\_sensitivity()}\spxextra{SimulationResult method}}

\begin{fulllineitems}
\phantomsection\label{\detokenize{results:SimulationResult.get_sensitivity}}
\pysigstartsignatures
\pysiglinewithargsret
{\sphinxbfcode{\sphinxupquote{get\_sensitivity}}}
{\sphinxparam{\DUrole{n}{node}}\sphinxparamcomma \sphinxparam{\DUrole{n}{param}}}
{}
\pysigstopsignatures
\sphinxAtStartPar
Returns the sensitivity gradient for a given output node and parameter.
For single‑point analyses it returns a scalar; otherwise a 1‑D array
across the sweep axis.

\end{fulllineitems}


\end{fulllineitems}


\sphinxstepscope


\chapter{Time‑Domain Waveform Module}
\label{\detokenize{waveforms:timedomain-waveform-module}}\label{\detokenize{waveforms:waveforms-module}}\label{\detokenize{waveforms::doc}}
\sphinxAtStartPar
This module implements a small library of transient source generators that are used during
transient analysis.  It supports the common SPICE waveforms \sphinxstylestrong{PULSE}, \sphinxstylestrong{SIN/SINE/COS} and \sphinxstylestrong{PWL}
( Piece‑Wise Linear ).  The {\hyperref[\detokenize{waveforms:Waveform}]{\sphinxcrossref{\sphinxcode{\sphinxupquote{Waveform}}}}} class parses the user‑supplied dictionary,
stores the waveform type, and evaluates the instantaneous voltage or current at any time \sphinxtitleref{t}.


\bigskip\hrule\bigskip

\index{Waveform (built\sphinxhyphen{}in class)@\spxentry{Waveform}\spxextra{built\sphinxhyphen{}in class}}

\begin{fulllineitems}
\phantomsection\label{\detokenize{waveforms:Waveform}}
\pysigstartsignatures
\pysiglinewithargsret
{\sphinxbfcode{\sphinxupquote{\DUrole{k}{class}\DUrole{w}{ }}}\sphinxbfcode{\sphinxupquote{Waveform}}}
{\sphinxparam{\DUrole{n}{source\_dict}}}
{}
\pysigstopsignatures
\sphinxAtStartPar
Evaluates time‑domain source functions for transient simulations.
\index{type (Waveform attribute)@\spxentry{type}\spxextra{Waveform attribute}}

\begin{fulllineitems}
\phantomsection\label{\detokenize{waveforms:Waveform.type}}
\pysigstartsignatures
\pysigline
{\sphinxbfcode{\sphinxupquote{type}}}
\pysigstopsignatures
\sphinxAtStartPar
\sphinxcode{\sphinxupquote{str}} \textendash{} The waveform type identifier (e.g. \sphinxcode{\sphinxupquote{\textquotesingle{}PULSE\textquotesingle{}}}, \sphinxcode{\sphinxupquote{\textquotesingle{}SIN\textquotesingle{}}}, \sphinxcode{\sphinxupquote{\textquotesingle{}COS\textquotesingle{}}}, \sphinxcode{\sphinxupquote{\textquotesingle{}PWL\textquotesingle{}}}).

\end{fulllineitems}

\index{params (Waveform attribute)@\spxentry{params}\spxextra{Waveform attribute}}

\begin{fulllineitems}
\phantomsection\label{\detokenize{waveforms:Waveform.params}}
\pysigstartsignatures
\pysigline
{\sphinxbfcode{\sphinxupquote{params}}}
\pysigstopsignatures
\sphinxAtStartPar
\sphinxcode{\sphinxupquote{dict}} \textendash{} All parameters that define the waveform shape, as parsed from a SPICE netlist.

\end{fulllineitems}

\index{get\_value() (Waveform method)@\spxentry{get\_value()}\spxextra{Waveform method}}

\begin{fulllineitems}
\phantomsection\label{\detokenize{waveforms:Waveform.get_value}}
\pysigstartsignatures
\pysiglinewithargsret
{\sphinxbfcode{\sphinxupquote{get\_value}}}
{\sphinxparam{\DUrole{n}{t}}}
{}
\pysigstopsignatures
\sphinxAtStartPar
Computes the instantaneous value of the source at simulation time \sphinxtitleref{t}.
\begin{quote}\begin{description}
\sphinxlineitem{Parameters}
\sphinxAtStartPar
\sphinxstyleliteralstrong{\sphinxupquote{t}} (\sphinxstyleliteralemphasis{\sphinxupquote{float}}) \textendash{} Current simulation time in seconds.

\sphinxlineitem{Returns}
\sphinxAtStartPar
\sphinxcode{\sphinxupquote{float}} \textendash{} Instantaneous voltage (V) or current (A).

\end{description}\end{quote}

\sphinxAtStartPar
The implementation supports the following waveform types:
\begin{itemize}
\item {} 
\sphinxAtStartPar
\sphinxstylestrong{PULSE}
Parameters: \sphinxcode{\sphinxupquote{V1, V2, TD, TR, TF, PW, PER}}
Returns a piecewise linear pulse with optional rise/fall times.

\item {} 
\sphinxAtStartPar
\sphinxstylestrong{SIN / SINE / COS}
Parameters: \sphinxcode{\sphinxupquote{VOFF, VAMP, FREQ, PHASE}} (phase is supplied in degrees).
Returns either a sine or cosine waveform depending on the type string.

\item {} 
\sphinxAtStartPar
\sphinxstylestrong{PWL}
Parameter: \sphinxcode{\sphinxupquote{TIME\_VOLTAGE\_PAIRS}} \textendash{} list of \sphinxcode{\sphinxupquote{(time, voltage)}} tuples.
Uses \sphinxcode{\sphinxupquote{numpy.interp()}} for linear interpolation between points.

\item {} 
\sphinxAtStartPar
\sphinxstylestrong{Static / Unknown} \textendash{} falls back to the value stored in \sphinxcode{\sphinxupquote{params{[}\textquotesingle{}value\textquotesingle{}{]}}} (defaults to 0).

\end{itemize}

\end{fulllineitems}


\end{fulllineitems}



\bigskip\hrule\bigskip



\section{Example usage}
\label{\detokenize{waveforms:example-usage}}
\begin{sphinxVerbatim}[commandchars=\\\{\}]
\PYG{k+kn}{from}\PYG{+w}{ }\PYG{n+nn}{waveforms}\PYG{+w}{ }\PYG{k+kn}{import} \PYG{n}{Waveform}

\PYG{c+c1}{\PYGZsh{} PULSE example}
\PYG{n}{pulse} \PYG{o}{=} \PYG{n}{Waveform}\PYG{p}{(}\PYG{p}{\PYGZob{}}
    \PYG{l+s+s2}{\PYGZdq{}}\PYG{l+s+s2}{type}\PYG{l+s+s2}{\PYGZdq{}}\PYG{p}{:} \PYG{l+s+s2}{\PYGZdq{}}\PYG{l+s+s2}{PULSE}\PYG{l+s+s2}{\PYGZdq{}}\PYG{p}{,}
    \PYG{l+s+s2}{\PYGZdq{}}\PYG{l+s+s2}{V1}\PYG{l+s+s2}{\PYGZdq{}}\PYG{p}{:} \PYG{l+m+mf}{0.0}\PYG{p}{,}
    \PYG{l+s+s2}{\PYGZdq{}}\PYG{l+s+s2}{V2}\PYG{l+s+s2}{\PYGZdq{}}\PYG{p}{:} \PYG{l+m+mf}{5.0}\PYG{p}{,}
    \PYG{l+s+s2}{\PYGZdq{}}\PYG{l+s+s2}{TD}\PYG{l+s+s2}{\PYGZdq{}}\PYG{p}{:} \PYG{l+m+mf}{1e\PYGZhy{}6}\PYG{p}{,}
    \PYG{l+s+s2}{\PYGZdq{}}\PYG{l+s+s2}{TR}\PYG{l+s+s2}{\PYGZdq{}}\PYG{p}{:} \PYG{l+m+mf}{10e\PYGZhy{}9}\PYG{p}{,}
    \PYG{l+s+s2}{\PYGZdq{}}\PYG{l+s+s2}{TF}\PYG{l+s+s2}{\PYGZdq{}}\PYG{p}{:} \PYG{l+m+mf}{10e\PYGZhy{}9}\PYG{p}{,}
    \PYG{l+s+s2}{\PYGZdq{}}\PYG{l+s+s2}{PW}\PYG{l+s+s2}{\PYGZdq{}}\PYG{p}{:} \PYG{l+m+mf}{100e\PYGZhy{}9}\PYG{p}{,}
    \PYG{l+s+s2}{\PYGZdq{}}\PYG{l+s+s2}{PER}\PYG{l+s+s2}{\PYGZdq{}}\PYG{p}{:} \PYG{l+m+mf}{200e\PYGZhy{}9}
\PYG{p}{\PYGZcb{}}\PYG{p}{)}

\PYG{n}{voltage\PYGZus{}at\PYGZus{}t} \PYG{o}{=} \PYG{n}{pulse}\PYG{o}{.}\PYG{n}{get\PYGZus{}value}\PYG{p}{(}\PYG{l+m+mf}{5e\PYGZhy{}7}\PYG{p}{)}  \PYG{c+c1}{\PYGZsh{} \PYGZhy{}\PYGZgt{} 5.0 V}
\end{sphinxVerbatim}

\sphinxstepscope


\chapter{AC Analysis Engine}
\label{\detokenize{ac_engine:ac-analysis-engine}}\label{\detokenize{ac_engine:ac-engine-module}}\label{\detokenize{ac_engine::doc}}
\sphinxAtStartPar
This module implements {\hyperref[\detokenize{ac_engine:ACEngine}]{\sphinxcrossref{\sphinxcode{\sphinxupquote{ACEngine}}}}}, which performs small‑signal frequency domain simulations.
The routine linearises all nonlinear components around a supplied DC operating point,
then sweeps the specified frequency range to compute complex phasors at every node.


\bigskip\hrule\bigskip

\index{ACEngine (built\sphinxhyphen{}in class)@\spxentry{ACEngine}\spxextra{built\sphinxhyphen{}in class}}

\begin{fulllineitems}
\phantomsection\label{\detokenize{ac_engine:ACEngine}}
\pysigstartsignatures
\pysiglinewithargsret
{\sphinxbfcode{\sphinxupquote{\DUrole{k}{class}\DUrole{w}{ }}}\sphinxbfcode{\sphinxupquote{ACEngine}}}
{\sphinxparam{\DUrole{n}{circuit}}\sphinxparamcomma \sphinxparam{\DUrole{n}{is\_nonlinear}}}
{}
\pysigstopsignatures
\sphinxAtStartPar
Core class that runs AC sweeps.
\index{circuit (ACEngine attribute)@\spxentry{circuit}\spxextra{ACEngine attribute}}

\begin{fulllineitems}
\phantomsection\label{\detokenize{ac_engine:ACEngine.circuit}}
\pysigstartsignatures
\pysigline
{\sphinxbfcode{\sphinxupquote{circuit}}}
\pysigstopsignatures
\sphinxAtStartPar
\sphinxcode{\sphinxupquote{Circuit}} \textendash{} The fully populated circuit object.

\end{fulllineitems}

\index{is\_nonlinear (ACEngine attribute)@\spxentry{is\_nonlinear}\spxextra{ACEngine attribute}}

\begin{fulllineitems}
\phantomsection\label{\detokenize{ac_engine:ACEngine.is_nonlinear}}
\pysigstartsignatures
\pysigline
{\sphinxbfcode{\sphinxupquote{is\_nonlinear}}}
\pysigstopsignatures
\sphinxAtStartPar
\sphinxcode{\sphinxupquote{bool}} \textendash{} True if the circuit contains nonlinear devices (diodes,
MOSFETs, etc.) that must be linearised at each bias point.

\end{fulllineitems}


\end{fulllineitems}



\bigskip\hrule\bigskip

\index{\_solve\_single\_point() (ACEngine method)@\spxentry{\_solve\_single\_point()}\spxextra{ACEngine method}}

\begin{fulllineitems}
\phantomsection\label{\detokenize{ac_engine:ACEngine._solve_single_point}}
\pysigstartsignatures
\pysiglinewithargsret
{\sphinxcode{\sphinxupquote{ACEngine.}}\sphinxbfcode{\sphinxupquote{\_solve\_single\_point}}}
{\sphinxparam{\DUrole{n}{w}}\sphinxparamcomma \sphinxparam{\DUrole{n}{Y\_small\_signal\_base}}}
{}
\pysigstopsignatures
\sphinxAtStartPar
Internal helper that solves the AC circuit for a single angular frequency.


\section{Parameters}
\label{\detokenize{ac_engine:parameters}}\begin{description}
\sphinxlineitem{w}{[}float{]}
\sphinxAtStartPar
Angular frequency (rad/s) = 2π·f.

\sphinxlineitem{Y\_small\_signal\_base}{[}\sphinxcode{\sphinxupquote{scipy.sparse.lil\_matrix}}{]}
\sphinxAtStartPar
Pre‑linearised admittance matrix that already contains static resistors
and the linearised conductances of any nonlinear devices.

\end{description}


\section{Returns}
\label{\detokenize{ac_engine:returns}}\begin{description}
\sphinxlineitem{tuple}
\sphinxAtStartPar
\sphinxcode{\sphinxupquote{(lu\_factorization, VI\_complex\_array)}} \textendash{} LU factorisation object and the
complex solution vector at this frequency.

\end{description}

\end{fulllineitems}



\bigskip\hrule\bigskip



\begin{fulllineitems}

\pysigstartsignatures
\pysigline
{\sphinxbfcode{\sphinxupquote{ACEngine.run(Y\_base\_lil,~VI\_dc,~start\_freq,~stop\_freq,~points,}}}
\pysigline
{\sphinxbfcode{\sphinxupquote{sweep\_type="DEC",~keep\_lus=False)}}}
\pysigstopsignatures
\sphinxAtStartPar
Executes an AC frequency sweep.


\section{Parameters}
\label{\detokenize{ac_engine:id1}}\begin{description}
\sphinxlineitem{Y\_base\_lil}{[}\sphinxcode{\sphinxupquote{scipy.sparse.lil\_matrix}}{]}
\sphinxAtStartPar
Static base admittance matrix (time‑invariant part of the circuit).

\sphinxlineitem{VI\_dc}{[}\sphinxcode{\sphinxupquote{numpy.ndarray}}{]}
\sphinxAtStartPar
DC operating point vector that was previously computed.

\sphinxlineitem{start\_freq, stop\_freq}{[}float{]}
\sphinxAtStartPar
Sweep bounds in hertz.

\sphinxlineitem{points}{[}int{]}
\sphinxAtStartPar
Number of frequency samples to generate.

\sphinxlineitem{sweep\_type}{[}str, optional{]}
\sphinxAtStartPar
\sphinxcode{\sphinxupquote{"DEC"}} (logarithmic), \sphinxcode{\sphinxupquote{"LIN"}}, \sphinxcode{\sphinxupquote{"OCT"}} or \sphinxcode{\sphinxupquote{"LIST"}}.  Default is \sphinxcode{\sphinxupquote{"DEC"}}.

\sphinxlineitem{keep\_lus}{[}bool, optional{]}
\sphinxAtStartPar
If True, caches the LU factorisation for every frequency point \textendash{} required by
the adjoint sensitivity engine.

\end{description}


\section{Returns}
\label{\detokenize{ac_engine:id2}}\begin{description}
\sphinxlineitem{tuple}
\sphinxAtStartPar
\sphinxcode{\sphinxupquote{(frequencies, VIs, list\_of\_lus)}}
\begin{itemize}
\item {} 
\sphinxAtStartPar
\sphinxstylestrong{frequencies} \textendash{} 1‑D array of sweep frequencies.

\item {} 
\sphinxAtStartPar
\sphinxstylestrong{VIs} \textendash{} 2‑D array (n\_frequencies × n\_nodes) containing complex phasors.

\item {} 
\sphinxAtStartPar
\sphinxstylestrong{list\_of\_lus} \textendash{} Optional list of LU objects, one per frequency.

\end{itemize}

\end{description}


\section{Notes}
\label{\detokenize{ac_engine:notes}}\begin{itemize}
\item {} 
\sphinxAtStartPar
The routine first “bakes” the small‑signal conductances into a base matrix
(\sphinxtitleref{Y\_small\_signal\_base}) once.  This optimisation avoids re‑stamping at every
frequency step.

\item {} 
\sphinxAtStartPar
For nonlinear circuits it calls each component’s
\sphinxcode{\sphinxupquote{stamp\_nonlinear()}} with the DC bias vector.

\end{itemize}

\end{fulllineitems}



\bigskip\hrule\bigskip



\section{Example usage}
\label{\detokenize{ac_engine:example-usage}}
\begin{sphinxVerbatim}[commandchars=\\\{\}]
\PYG{k+kn}{from}\PYG{+w}{ }\PYG{n+nn}{core}\PYG{n+nn}{.}\PYG{n+nn}{circuit}\PYG{+w}{ }\PYG{k+kn}{import} \PYG{n}{Circuit}
\PYG{k+kn}{from}\PYG{+w}{ }\PYG{n+nn}{ac\PYGZus{}engine}\PYG{+w}{ }\PYG{k+kn}{import} \PYG{n}{ACEngine}

\PYG{n}{engine} \PYG{o}{=} \PYG{n}{ACEngine}\PYG{p}{(}\PYG{n}{circuit}\PYG{p}{,} \PYG{n}{is\PYGZus{}nonlinear}\PYG{o}{=}\PYG{k+kc}{True}\PYG{p}{)}
\PYG{n}{freqs}\PYG{p}{,} \PYG{n}{results}\PYG{p}{,} \PYG{n}{luts} \PYG{o}{=} \PYG{n}{engine}\PYG{o}{.}\PYG{n}{run}\PYG{p}{(}
    \PYG{n}{Y\PYGZus{}base\PYGZus{}lil}\PYG{o}{=}\PYG{n}{Y\PYGZus{}base}\PYG{p}{,}
    \PYG{n}{VI\PYGZus{}dc}\PYG{o}{=}\PYG{n}{dc\PYGZus{}solution}\PYG{p}{,}
    \PYG{n}{start\PYGZus{}freq}\PYG{o}{=}\PYG{l+m+mf}{1e3}\PYG{p}{,}
    \PYG{n}{stop\PYGZus{}freq}\PYG{o}{=}\PYG{l+m+mf}{1e6}\PYG{p}{,}
    \PYG{n}{points}\PYG{o}{=}\PYG{l+m+mi}{100}
\PYG{p}{)}
\end{sphinxVerbatim}

\sphinxstepscope


\chapter{DC Analysis Engine}
\label{\detokenize{dc_engine:dc-analysis-engine}}\label{\detokenize{dc_engine:dc-engine-module}}\label{\detokenize{dc_engine::doc}}
\sphinxAtStartPar
This module implements {\hyperref[\detokenize{dc_engine:DCEngine}]{\sphinxcrossref{\sphinxcode{\sphinxupquote{DCEngine}}}}}, which is responsible for building the static
Modified Nodal Analysis (MNA) topology and solving the DC operating point.
The resulting bias solution is required before any AC or transient analysis can begin.


\bigskip\hrule\bigskip

\index{DCEngine (built\sphinxhyphen{}in class)@\spxentry{DCEngine}\spxextra{built\sphinxhyphen{}in class}}

\begin{fulllineitems}
\phantomsection\label{\detokenize{dc_engine:DCEngine}}
\pysigstartsignatures
\pysiglinewithargsret
{\sphinxbfcode{\sphinxupquote{\DUrole{k}{class}\DUrole{w}{ }}}\sphinxbfcode{\sphinxupquote{DCEngine}}}
{\sphinxparam{\DUrole{n}{circuit}}\sphinxparamcomma \sphinxparam{\DUrole{n}{is\_complex}}\sphinxparamcomma \sphinxparam{\DUrole{n}{is\_nonlinear}}}
{}
\pysigstopsignatures
\sphinxAtStartPar
Core class that compiles the circuit’s base matrices and performs
steady‑state solves.
\index{circuit (DCEngine attribute)@\spxentry{circuit}\spxextra{DCEngine attribute}}

\begin{fulllineitems}
\phantomsection\label{\detokenize{dc_engine:DCEngine.circuit}}
\pysigstartsignatures
\pysigline
{\sphinxbfcode{\sphinxupquote{circuit}}}
\pysigstopsignatures
\sphinxAtStartPar
\sphinxcode{\sphinxupquote{Circuit}} \textendash{} The fully populated circuit object.

\end{fulllineitems}

\index{is\_complex (DCEngine attribute)@\spxentry{is\_complex}\spxextra{DCEngine attribute}}

\begin{fulllineitems}
\phantomsection\label{\detokenize{dc_engine:DCEngine.is_complex}}
\pysigstartsignatures
\pysigline
{\sphinxbfcode{\sphinxupquote{is\_complex}}}
\pysigstopsignatures
\sphinxAtStartPar
\sphinxcode{\sphinxupquote{bool}} \textendash{} True if an AC analysis has been requested; forces complex
matrix data types.

\end{fulllineitems}

\index{is\_nonlinear (DCEngine attribute)@\spxentry{is\_nonlinear}\spxextra{DCEngine attribute}}

\begin{fulllineitems}
\phantomsection\label{\detokenize{dc_engine:DCEngine.is_nonlinear}}
\pysigstartsignatures
\pysigline
{\sphinxbfcode{\sphinxupquote{is\_nonlinear}}}
\pysigstopsignatures
\sphinxAtStartPar
\sphinxcode{\sphinxupquote{bool}} \textendash{} True if any component in the circuit requires a Newton‑Raphson
iteration (e.g. MOSFETs, diodes).

\end{fulllineitems}


\end{fulllineitems}



\bigskip\hrule\bigskip



\bigskip\hrule\bigskip



\bigskip\hrule\bigskip



\bigskip\hrule\bigskip



\section{Example usage}
\label{\detokenize{dc_engine:example-usage}}
\begin{sphinxVerbatim}[commandchars=\\\{\}]
\PYG{k+kn}{from}\PYG{+w}{ }\PYG{n+nn}{core}\PYG{n+nn}{.}\PYG{n+nn}{circuit}\PYG{+w}{ }\PYG{k+kn}{import} \PYG{n}{Circuit}
\PYG{k+kn}{from}\PYG{+w}{ }\PYG{n+nn}{dc\PYGZus{}engine}\PYG{+w}{ }\PYG{k+kn}{import} \PYG{n}{DCEngine}

\PYG{c+c1}{\PYGZsh{} Assume `circuit` has been built and we want a DC sweep of V1}
\PYG{n}{engine} \PYG{o}{=} \PYG{n}{DCEngine}\PYG{p}{(}\PYG{n}{circuit}\PYG{p}{,} \PYG{n}{is\PYGZus{}complex}\PYG{o}{=}\PYG{k+kc}{False}\PYG{p}{,} \PYG{n}{is\PYGZus{}nonlinear}\PYG{o}{=}\PYG{k+kc}{True}\PYG{p}{)}
\PYG{n}{Y\PYGZus{}base}\PYG{p}{,} \PYG{n}{\PYGZus{}} \PYG{o}{=} \PYG{n}{engine}\PYG{o}{.}\PYG{n}{build\PYGZus{}base\PYGZus{}matrices}\PYG{p}{(}\PYG{p}{)}
\PYG{n}{axis}\PYG{p}{,} \PYG{n}{Vs}\PYG{p}{,} \PYG{n}{lus} \PYG{o}{=} \PYG{n}{engine}\PYG{o}{.}\PYG{n}{compute\PYGZus{}dc\PYGZus{}sweep}\PYG{p}{(}\PYG{n}{Y\PYGZus{}base}\PYG{p}{,} \PYG{l+s+s1}{\PYGZsq{}}\PYG{l+s+s1}{V1}\PYG{l+s+s1}{\PYGZsq{}}\PYG{p}{,} \PYG{l+m+mf}{0.0}\PYG{p}{,} \PYG{l+m+mf}{5.0}\PYG{p}{,} \PYG{l+m+mf}{0.5}\PYG{p}{)}
\end{sphinxVerbatim}

\sphinxstepscope


\chapter{Adjoint Sensitivity Engine}
\label{\detokenize{adjoint_engine:adjoint-sensitivity-engine}}\label{\detokenize{adjoint_engine:adjoint-engine-module}}\label{\detokenize{adjoint_engine::doc}}
\sphinxAtStartPar
This module implements {\hyperref[\detokenize{adjoint_engine:AdjointEngine}]{\sphinxcrossref{\sphinxcode{\sphinxupquote{AdjointEngine}}}}}, which evaluates parameter sensitivities using the adjoint network method.
The engine supports both \sphinxstyleemphasis{continuous} (steady‑state, AC, DC, OP) and \sphinxstyleemphasis{global transient} sensitivity calculations.


\bigskip\hrule\bigskip

\index{AdjointEngine (built\sphinxhyphen{}in class)@\spxentry{AdjointEngine}\spxextra{built\sphinxhyphen{}in class}}

\begin{fulllineitems}
\phantomsection\label{\detokenize{adjoint_engine:AdjointEngine}}
\pysigstartsignatures
\pysiglinewithargsret
{\sphinxbfcode{\sphinxupquote{\DUrole{k}{class}\DUrole{w}{ }}}\sphinxbfcode{\sphinxupquote{AdjointEngine}}}
{\sphinxparam{\DUrole{n}{circuit}}\sphinxparamcomma \sphinxparam{\DUrole{n}{output\_nodes}}}
{}
\pysigstopsignatures
\sphinxAtStartPar
Initializes the adjoint solver for a given circuit.
\index{circuit (AdjointEngine attribute)@\spxentry{circuit}\spxextra{AdjointEngine attribute}}

\begin{fulllineitems}
\phantomsection\label{\detokenize{adjoint_engine:AdjointEngine.circuit}}
\pysigstartsignatures
\pysigline
{\sphinxbfcode{\sphinxupquote{circuit}}}
\pysigstopsignatures
\sphinxAtStartPar
\sphinxcode{\sphinxupquote{Circuit}} \textendash{} The main circuit object containing components.

\end{fulllineitems}

\index{output\_nodes (AdjointEngine attribute)@\spxentry{output\_nodes}\spxextra{AdjointEngine attribute}}

\begin{fulllineitems}
\phantomsection\label{\detokenize{adjoint_engine:AdjointEngine.output_nodes}}
\pysigstartsignatures
\pysigline
{\sphinxbfcode{\sphinxupquote{output\_nodes}}}
\pysigstopsignatures
\sphinxAtStartPar
\sphinxcode{\sphinxupquote{list{[}str{]}}} \textendash{} Nodes that will be used as objectives in the sensitivity analysis.
The constructor resolves node names against \sphinxcode{\sphinxupquote{circuit.node\_map}} so that
strings, integers and mixed types all work correctly.

\end{fulllineitems}


\end{fulllineitems}



\bigskip\hrule\bigskip

\index{\_solve\_adjoint() (AdjointEngine method)@\spxentry{\_solve\_adjoint()}\spxextra{AdjointEngine method}}

\begin{fulllineitems}
\phantomsection\label{\detokenize{adjoint_engine:AdjointEngine._solve_adjoint}}
\pysigstartsignatures
\pysiglinewithargsret
{\sphinxcode{\sphinxupquote{AdjointEngine.}}\sphinxbfcode{\sphinxupquote{\_solve\_adjoint}}}
{\sphinxparam{\DUrole{n}{lu}}\sphinxparamcomma \sphinxparam{\DUrole{n}{target}}\sphinxparamcomma \sphinxparam{\DUrole{n}{base\_rhs}\DUrole{o}{=}\DUrole{default_value}{None}}\sphinxparamcomma \sphinxparam{\DUrole{n}{is\_complex}\DUrole{o}{=}\DUrole{default_value}{False}}}
{}
\pysigstopsignatures
\sphinxAtStartPar
Internal helper that solves the transposed matrix equation for a single target node.
The method can be reused by both transient and continuous solvers.


\section{Parameters}
\label{\detokenize{adjoint_engine:parameters}}\begin{description}
\sphinxlineitem{lu}{[}LU factorization object{]}
\sphinxAtStartPar
From the forward solver (e.g., SciPy or custom).

\sphinxlineitem{target}{[}str | None{]}
\sphinxAtStartPar
Node name that will receive an impulse of +1.0 in the RHS.

\sphinxlineitem{base\_rhs}{[}array\_like, optional{]}
\sphinxAtStartPar
Optional additional RHS vector to be added before solving.

\sphinxlineitem{is\_complex}{[}bool, optional{]}
\sphinxAtStartPar
Flag indicating whether complex arithmetic is required (AC analysis).

\end{description}


\section{Returns}
\label{\detokenize{adjoint_engine:returns}}\begin{description}
\sphinxlineitem{array\_like}
\sphinxAtStartPar
Solution vector \sphinxtitleref{ψ} (adjoint variable) with shape \sphinxtitleref{(total\_dim,)}.

\end{description}

\end{fulllineitems}



\bigskip\hrule\bigskip

\index{compute\_transient() (AdjointEngine method)@\spxentry{compute\_transient()}\spxextra{AdjointEngine method}}

\begin{fulllineitems}
\phantomsection\label{\detokenize{adjoint_engine:AdjointEngine.compute_transient}}
\pysigstartsignatures
\pysiglinewithargsret
{\sphinxcode{\sphinxupquote{AdjointEngine.}}\sphinxbfcode{\sphinxupquote{compute\_transient}}}
{\sphinxparam{\DUrole{n}{time\_array}}\sphinxparamcomma \sphinxparam{\DUrole{n}{V\_forward}}\sphinxparamcomma \sphinxparam{\DUrole{n}{list\_of\_lus}}\sphinxparamcomma \sphinxparam{\DUrole{n}{dt}}\sphinxparamcomma \sphinxparam{\DUrole{n}{method}\DUrole{o}{=}\DUrole{default_value}{\textquotesingle{}BE\textquotesingle{}}}}
{}
\pysigstopsignatures
\sphinxAtStartPar
Calculates exact \sphinxstyleemphasis{global} adjoint sensitivities via backward time integration.
The result is a dictionary containing:
\begin{itemize}
\item {} 
\sphinxAtStartPar
\sphinxstylestrong{Integrated\_Transient} \textendash{} scalar integrals for each parameter and target node.

\item {} 
\sphinxAtStartPar
\sphinxstylestrong{Time\_Series} \textendash{} a \sphinxcode{\sphinxupquote{SensitivityData}} tensor that stores the full time‑series of sensitivities.

\item {} 
\sphinxAtStartPar
\sphinxstylestrong{Raw\_Adjoint\_History} \textendash{} raw adjoint vectors at every step (useful for fault analysis).

\end{itemize}


\section{Parameters}
\label{\detokenize{adjoint_engine:id1}}\begin{description}
\sphinxlineitem{time\_array}{[}1‑D array\_like{]}
\sphinxAtStartPar
Simulation time points.

\sphinxlineitem{V\_forward}{[}list/array{]}
\sphinxAtStartPar
Forward solution vector(s) at each time step.

\sphinxlineitem{list\_of\_lus}{[}list of LU objects{]}
\sphinxAtStartPar
LU factorizations from the forward transient solver, one per time step.

\sphinxlineitem{dt}{[}float{]}
\sphinxAtStartPar
Time step size.

\sphinxlineitem{method}{[}str, optional{]}
\sphinxAtStartPar
Integration scheme used in the forward pass (\sphinxtitleref{“BE”} = Backward Euler,
\sphinxtitleref{“TR”} = Trapezoidal).  The same method is used for the backward solve.

\end{description}


\section{Notes}
\label{\detokenize{adjoint_engine:notes}}\begin{itemize}
\item {} 
\sphinxAtStartPar
The algorithm performs a reverse‑time sweep to compute adjoint variables
and then a forward integration to accumulate sensitivities.

\item {} 
\sphinxAtStartPar
All parameters that are differentiable in the circuit (from
\sphinxcode{\sphinxupquote{circuit.differentiable\_params}}) are included in the tensor.

\end{itemize}

\end{fulllineitems}



\bigskip\hrule\bigskip



\begin{fulllineitems}

\pysigstartsignatures
\pysigline
{\sphinxbfcode{\sphinxupquote{AdjointEngine.compute\_sensitivities(sweep\_axis,~V\_forward,~list\_of\_lus,}}}
\pysigline
{\sphinxbfcode{\sphinxupquote{domain="static",~method="TR",~dt=0.0)}}}
\pysigstopsignatures
\sphinxAtStartPar
Unified continuous sensitivity solver for AC, DC, OP and local‑DC (transient) analyses.
The routine builds a \sphinxcode{\sphinxupquote{SensitivityData}} tensor that
contains the partial derivatives of each output node with respect to every
differentiable parameter.


\section{Parameters}
\label{\detokenize{adjoint_engine:id2}}\begin{description}
\sphinxlineitem{sweep\_axis}{[}1‑D array\_like{]}
\sphinxAtStartPar
Independent variable array (time, frequency, voltage, or single value for OP).

\sphinxlineitem{V\_forward}{[}list/array{]}
\sphinxAtStartPar
Forward solution vectors at each sweep point.

\sphinxlineitem{list\_of\_lus}{[}list of LU objects{]}
\sphinxAtStartPar
LU factorizations from the forward solver, one per step.

\sphinxlineitem{domain}{[}str, optional{]}
\sphinxAtStartPar
One of \sphinxcode{\sphinxupquote{"time"}}, \sphinxcode{\sphinxupquote{"frequency"}}, \sphinxcode{\sphinxupquote{"voltage"}}, or \sphinxcode{\sphinxupquote{"static"}}.

\sphinxlineitem{method}{[}str, optional{]}
\sphinxAtStartPar
Integration scheme used for transient problems (\sphinxtitleref{“TR”} or \sphinxtitleref{“BE”}).

\sphinxlineitem{dt}{[}float, optional{]}
\sphinxAtStartPar
Time step size (used only when \sphinxcode{\sphinxupquote{domain == "time"}}).

\end{description}


\section{Returns}
\label{\detokenize{adjoint_engine:id3}}\begin{description}
\sphinxlineitem{SensitivityData}
\sphinxAtStartPar
Tensor containing all continuous sensitivities.  The object also stores the
raw adjoint vectors in its \sphinxcode{\sphinxupquote{adjoint\_vectors}} attribute for later fault analysis.

\end{description}

\end{fulllineitems}



\bigskip\hrule\bigskip



\section{Example usage}
\label{\detokenize{adjoint_engine:example-usage}}
\begin{sphinxVerbatim}[commandchars=\\\{\}]
\PYG{k+kn}{from}\PYG{+w}{ }\PYG{n+nn}{core}\PYG{n+nn}{.}\PYG{n+nn}{circuit}\PYG{+w}{ }\PYG{k+kn}{import} \PYG{n}{Circuit}
\PYG{k+kn}{from}\PYG{+w}{ }\PYG{n+nn}{adjoint\PYGZus{}engine}\PYG{+w}{ }\PYG{k+kn}{import} \PYG{n}{AdjointEngine}

\PYG{c+c1}{\PYGZsh{} Assume `circuit` has been built and a forward solution is available.}
\PYG{n}{engine} \PYG{o}{=} \PYG{n}{AdjointEngine}\PYG{p}{(}\PYG{n}{circuit}\PYG{p}{,} \PYG{n}{output\PYGZus{}nodes}\PYG{o}{=}\PYG{p}{[}\PYG{l+s+s1}{\PYGZsq{}}\PYG{l+s+s1}{out}\PYG{l+s+s1}{\PYGZsq{}}\PYG{p}{]}\PYG{p}{)}
\PYG{n}{tensor} \PYG{o}{=} \PYG{n}{engine}\PYG{o}{.}\PYG{n}{compute\PYGZus{}sensitivities}\PYG{p}{(}
    \PYG{n}{sweep\PYGZus{}axis}\PYG{o}{=}\PYG{n}{np}\PYG{o}{.}\PYG{n}{linspace}\PYG{p}{(}\PYG{l+m+mi}{0}\PYG{p}{,} \PYG{l+m+mf}{1e\PYGZhy{}3}\PYG{p}{,} \PYG{l+m+mi}{100}\PYG{p}{)}\PYG{p}{,}
    \PYG{n}{V\PYGZus{}forward}\PYG{o}{=}\PYG{n}{forward\PYGZus{}results}\PYG{p}{,}
    \PYG{n}{list\PYGZus{}of\PYGZus{}lus}\PYG{o}{=}\PYG{n}{luts}\PYG{p}{,}
    \PYG{n}{domain}\PYG{o}{=}\PYG{l+s+s1}{\PYGZsq{}}\PYG{l+s+s1}{time}\PYG{l+s+s1}{\PYGZsq{}}\PYG{p}{,}
    \PYG{n}{method}\PYG{o}{=}\PYG{l+s+s1}{\PYGZsq{}}\PYG{l+s+s1}{TR}\PYG{l+s+s1}{\PYGZsq{}}
\PYG{p}{)}
\end{sphinxVerbatim}

\sphinxstepscope


\chapter{Graphical User Interface}
\label{\detokenize{gui:graphical-user-interface}}\label{\detokenize{gui:gui-module}}\label{\detokenize{gui::doc}}
\sphinxAtStartPar
This module implements a \sphinxstylestrong{Tkinter}‑based front end for the Python SPICE simulator.
It provides an interactive editor for netlists, asynchronous simulation execution,
and embedded Matplotlib figures to visualize transient, AC and DC results.


\bigskip\hrule\bigskip

\index{CircuitSimulatorGUI (built\sphinxhyphen{}in class)@\spxentry{CircuitSimulatorGUI}\spxextra{built\sphinxhyphen{}in class}}

\begin{fulllineitems}
\phantomsection\label{\detokenize{gui:CircuitSimulatorGUI}}
\pysigstartsignatures
\pysiglinewithargsret
{\sphinxbfcode{\sphinxupquote{\DUrole{k}{class}\DUrole{w}{ }}}\sphinxbfcode{\sphinxupquote{CircuitSimulatorGUI}}}
{\sphinxparam{\DUrole{n}{root}}\sphinxparamcomma \sphinxparam{\DUrole{n}{simulation\_callback}}}
{}
\pysigstopsignatures
\sphinxAtStartPar
Main application window that coordinates all GUI widgets and
communicates with the backend simulation engine.
\index{root (CircuitSimulatorGUI attribute)@\spxentry{root}\spxextra{CircuitSimulatorGUI attribute}}

\begin{fulllineitems}
\phantomsection\label{\detokenize{gui:CircuitSimulatorGUI.root}}
\pysigstartsignatures
\pysigline
{\sphinxbfcode{\sphinxupquote{root}}}
\pysigstopsignatures
\sphinxAtStartPar
\sphinxcode{\sphinxupquote{tk.Tk}} \textendash{} The main Tkinter window instance.

\end{fulllineitems}

\index{run\_simulation\_core (CircuitSimulatorGUI attribute)@\spxentry{run\_simulation\_core}\spxextra{CircuitSimulatorGUI attribute}}

\begin{fulllineitems}
\phantomsection\label{\detokenize{gui:CircuitSimulatorGUI.run_simulation_core}}
\pysigstartsignatures
\pysigline
{\sphinxbfcode{\sphinxupquote{run\_simulation\_core}}}
\pysigstopsignatures
\sphinxAtStartPar
Callable \textendash{} Reference to the \sphinxtitleref{run\_simulation\_core} function defined in
\sphinxcode{\sphinxupquote{main.py}}.  Invoked when the user clicks \sphinxstyleemphasis{Run Simulation}.

\end{fulllineitems}

\index{circuit (CircuitSimulatorGUI attribute)@\spxentry{circuit}\spxextra{CircuitSimulatorGUI attribute}}

\begin{fulllineitems}
\phantomsection\label{\detokenize{gui:CircuitSimulatorGUI.circuit}}
\pysigstartsignatures
\pysigline
{\sphinxbfcode{\sphinxupquote{circuit}}}
\pysigstopsignatures
\sphinxAtStartPar
\sphinxcode{\sphinxupquote{Circuit | None}} \textendash{} The parsed circuit object returned
by the backend after a simulation finishes.

\end{fulllineitems}

\index{result (CircuitSimulatorGUI attribute)@\spxentry{result}\spxextra{CircuitSimulatorGUI attribute}}

\begin{fulllineitems}
\phantomsection\label{\detokenize{gui:CircuitSimulatorGUI.result}}
\pysigstartsignatures
\pysigline
{\sphinxbfcode{\sphinxupquote{result}}}
\pysigstopsignatures
\sphinxAtStartPar
\sphinxcode{\sphinxupquote{SimulationResult | None}} \textendash{} The data vault holding all
forward and sensitivity results.

\end{fulllineitems}


\end{fulllineitems}



\bigskip\hrule\bigskip



\section{Key Methods}
\label{\detokenize{gui:key-methods}}

\bigskip\hrule\bigskip



\section{Additional Details}
\label{\detokenize{gui:additional-details}}\begin{itemize}
\item {} 
\sphinxAtStartPar
\sphinxstylestrong{Threading} \textendash{} Simulation runs on a daemon thread (\sphinxtitleref{threading.Thread}) to keep
the UI responsive.  All GUI updates are marshalled back to the main thread via
\sphinxtitleref{root.after}.

\item {} 
\sphinxAtStartPar
\sphinxstylestrong{Plotting} \textendash{} The widget layout uses a split pane: left side is a
read‑only text area for the netlist, right side is a notebook with two tabs:
\sphinxstyleemphasis{Waveform Plot} (Matplotlib) and \sphinxstyleemphasis{Nodal Analysis Data} (plain text).

\item {} 
\sphinxAtStartPar
\sphinxstylestrong{Node Naming} \textendash{} Nodes are displayed as \sphinxtitleref{V(node)} or \sphinxtitleref{I(V1)} depending on whether
they represent voltage nodes or MNA branch currents.  Voltage nodes are auto‑selected.

\item {} 
\sphinxAtStartPar
\sphinxstylestrong{Sensitivity UI} \textendash{} A combobox lists all parameters that the adjoint engine has
identified for each node; selecting one overlays its sensitivity curve.

\end{itemize}


\bigskip\hrule\bigskip


\sphinxAtStartPar
Example Usage (from a main script)

\begin{sphinxVerbatim}[commandchars=\\\{\}]
\PYG{k+kn}{import}\PYG{+w}{ }\PYG{n+nn}{tkinter}\PYG{+w}{ }\PYG{k}{as}\PYG{+w}{ }\PYG{n+nn}{tk}
\PYG{k+kn}{from}\PYG{+w}{ }\PYG{n+nn}{gui}\PYG{+w}{ }\PYG{k+kn}{import} \PYG{n}{CircuitSimulatorGUI}
\PYG{k+kn}{from}\PYG{+w}{ }\PYG{n+nn}{main}\PYG{+w}{ }\PYG{k+kn}{import} \PYG{n}{run\PYGZus{}simulation\PYGZus{}core}

\PYG{n}{root} \PYG{o}{=} \PYG{n}{tk}\PYG{o}{.}\PYG{n}{Tk}\PYG{p}{(}\PYG{p}{)}
\PYG{n}{app} \PYG{o}{=} \PYG{n}{CircuitSimulatorGUI}\PYG{p}{(}\PYG{n}{root}\PYG{p}{,} \PYG{n}{run\PYGZus{}simulation\PYGZus{}core}\PYG{p}{)}
\PYG{n}{root}\PYG{o}{.}\PYG{n}{mainloop}\PYG{p}{(}\PYG{p}{)}
\end{sphinxVerbatim}

\sphinxstepscope


\chapter{SPICE Netlist Parser}
\label{\detokenize{parser:spice-netlist-parser}}\label{\detokenize{parser:netlist-parser-module}}\label{\detokenize{parser::doc}}
\sphinxAtStartPar
This module reads standard SPICE‑formatted text files, handles line continuations,
extracts scaling suffixes (k, MEG, u, etc.) and compiles the data into dictionaries that
the \sphinxcode{\sphinxupquote{Circuit}} factory can consume.


\bigskip\hrule\bigskip

\index{NetlistParser (built\sphinxhyphen{}in class)@\spxentry{NetlistParser}\spxextra{built\sphinxhyphen{}in class}}

\begin{fulllineitems}
\phantomsection\label{\detokenize{parser:NetlistParser}}
\pysigstartsignatures
\pysigline
{\sphinxbfcode{\sphinxupquote{\DUrole{k}{class}\DUrole{w}{ }}}\sphinxbfcode{\sphinxupquote{NetlistParser}}}
\pysigstopsignatures
\sphinxAtStartPar
Core parser for SPICE netlists.
\index{components (NetlistParser attribute)@\spxentry{components}\spxextra{NetlistParser attribute}}

\begin{fulllineitems}
\phantomsection\label{\detokenize{parser:NetlistParser.components}}
\pysigstartsignatures
\pysigline
{\sphinxbfcode{\sphinxupquote{components}}}
\pysigstopsignatures
\sphinxAtStartPar
\sphinxcode{\sphinxupquote{dict}} \textendash{} Parsed component data keyed by component name (e.g. \sphinxcode{\sphinxupquote{\textquotesingle{}R1\textquotesingle{}}}).

\end{fulllineitems}

\index{models (NetlistParser attribute)@\spxentry{models}\spxextra{NetlistParser attribute}}

\begin{fulllineitems}
\phantomsection\label{\detokenize{parser:NetlistParser.models}}
\pysigstartsignatures
\pysigline
{\sphinxbfcode{\sphinxupquote{models}}}
\pysigstopsignatures
\sphinxAtStartPar
\sphinxcode{\sphinxupquote{dict}} \textendash{} Parsed \sphinxcode{\sphinxupquote{.MODEL}} definitions.

\end{fulllineitems}

\index{analyses (NetlistParser attribute)@\spxentry{analyses}\spxextra{NetlistParser attribute}}

\begin{fulllineitems}
\phantomsection\label{\detokenize{parser:NetlistParser.analyses}}
\pysigstartsignatures
\pysigline
{\sphinxbfcode{\sphinxupquote{analyses}}}
\pysigstopsignatures
\sphinxAtStartPar
\sphinxcode{\sphinxupquote{dict}} \textendash{} Parsed simulation commands (.TRAN, .AC, .DC, .OP).

\end{fulllineitems}

\index{current\_line (NetlistParser attribute)@\spxentry{current\_line}\spxextra{NetlistParser attribute}}

\begin{fulllineitems}
\phantomsection\label{\detokenize{parser:NetlistParser.current_line}}
\pysigstartsignatures
\pysigline
{\sphinxbfcode{\sphinxupquote{current\_line}}}
\pysigstopsignatures
\sphinxAtStartPar
\sphinxcode{\sphinxupquote{int}} \textendash{} Line number of the line currently being parsed (for error reporting).

\end{fulllineitems}

\index{raw\_line\_text (NetlistParser attribute)@\spxentry{raw\_line\_text}\spxextra{NetlistParser attribute}}

\begin{fulllineitems}
\phantomsection\label{\detokenize{parser:NetlistParser.raw_line_text}}
\pysigstartsignatures
\pysigline
{\sphinxbfcode{\sphinxupquote{raw\_line\_text}}}
\pysigstopsignatures
\sphinxAtStartPar
\sphinxcode{\sphinxupquote{str}} \textendash{} Raw text of the current logical line.

\end{fulllineitems}

\index{dispatch\_registry (NetlistParser attribute)@\spxentry{dispatch\_registry}\spxextra{NetlistParser attribute}}

\begin{fulllineitems}
\phantomsection\label{\detokenize{parser:NetlistParser.dispatch_registry}}
\pysigstartsignatures
\pysigline
{\sphinxbfcode{\sphinxupquote{dispatch\_registry}}}
\pysigstopsignatures
\sphinxAtStartPar
\sphinxcode{\sphinxupquote{dict}} \textendash{} Maps SPICE prefix letters to the corresponding mini‑parser
methods used by \sphinxcode{\sphinxupquote{\_parse\_component()}}.

\end{fulllineitems}


\end{fulllineitems}



\bigskip\hrule\bigskip



\section{Key Methods}
\label{\detokenize{parser:key-methods}}

\bigskip\hrule\bigskip



\section{Component Mini‑Parsers}
\label{\detokenize{parser:component-miniparsers}}\begin{itemize}
\item {} 
\sphinxAtStartPar
\sphinxcode{\sphinxupquote{\_parse\_passive}} \textendash{} Resistors, capacitors, inductors.
Returns dict with keys: \sphinxcode{\sphinxupquote{type}}, \sphinxcode{\sphinxupquote{n1}}, \sphinxcode{\sphinxupquote{n2}}, \sphinxcode{\sphinxupquote{value}}.

\item {} 
\sphinxAtStartPar
\sphinxcode{\sphinxupquote{\_parse\_diode}} \textendash{} Diodes; may include a model name or saturation current.

\item {} 
\sphinxAtStartPar
\sphinxcode{\sphinxupquote{\_parse\_mosfet}} \textendash{} MOSFETs (four‑terminal); parses node list, model and
instance parameters.

\item {} 
\sphinxAtStartPar
\sphinxcode{\sphinxupquote{\_parse\_source}} \textendash{} Voltage/Current sources; handles DC, AC and transient
definitions (PULSE, SIN, PWL etc.).

\item {} 
\sphinxAtStartPar
\sphinxcode{\sphinxupquote{\_parse\_vccs}}, \sphinxcode{\sphinxupquote{\_parse\_vcvs}}, \sphinxcode{\sphinxupquote{\_parse\_cccs}}, \sphinxcode{\sphinxupquote{\_parse\_ccvs}} \textendash{} Various
controlled sources.

\end{itemize}


\bigskip\hrule\bigskip



\section{Utility Functions}
\label{\detokenize{parser:utility-functions}}\index{\_parse\_node() (NetlistParser static method)@\spxentry{\_parse\_node()}\spxextra{NetlistParser static method}}

\begin{fulllineitems}
\phantomsection\label{\detokenize{parser:NetlistParser._parse_node}}
\pysigstartsignatures
\pysiglinewithargsret
{\sphinxbfcode{\sphinxupquote{\DUrole{k}{static}\DUrole{w}{ }}}\sphinxcode{\sphinxupquote{NetlistParser.}}\sphinxbfcode{\sphinxupquote{\_parse\_node}}}
{\sphinxparam{\DUrole{n}{node\_str}}}
{}
\pysigstopsignatures
\sphinxAtStartPar
Converts node identifiers to integers or keeps them as strings.
Recognises the ground nodes \sphinxcode{\sphinxupquote{0}} and \sphinxcode{\sphinxupquote{GND}}.

\end{fulllineitems}

\index{\_parse\_value() (NetlistParser static method)@\spxentry{\_parse\_value()}\spxextra{NetlistParser static method}}

\begin{fulllineitems}
\phantomsection\label{\detokenize{parser:NetlistParser._parse_value}}
\pysigstartsignatures
\pysiglinewithargsret
{\sphinxbfcode{\sphinxupquote{\DUrole{k}{static}\DUrole{w}{ }}}\sphinxcode{\sphinxupquote{NetlistParser.}}\sphinxbfcode{\sphinxupquote{\_parse\_value}}}
{\sphinxparam{\DUrole{n}{value\_str}}}
{}
\pysigstopsignatures
\sphinxAtStartPar
Converts SPICE numbers with scaling suffixes (e.g. \sphinxcode{\sphinxupquote{10k}}, \sphinxcode{\sphinxupquote{5u}}) into floats.
Supports multipliers: T, G, MEG, K, MIL, M, U, N, P, F.

\end{fulllineitems}



\bigskip\hrule\bigskip



\section{Helper Methods}
\label{\detokenize{parser:helper-methods}}\begin{itemize}
\item {} 
\sphinxAtStartPar
\sphinxcode{\sphinxupquote{\_preprocess\_lines(lines)}} \textendash{} Handles line continuations starting with \sphinxcode{\sphinxupquote{+}},
strips comments, and yields tuples \sphinxcode{\sphinxupquote{(line\_number, logical\_line)}}.

\item {} 
\sphinxAtStartPar
\sphinxcode{\sphinxupquote{\_attach\_models()}} \textendash{} Binds parsed \sphinxcode{\sphinxupquote{.MODEL}} definitions to the relevant
components, upgrading their type and adding \sphinxcode{\sphinxupquote{model\_params}}.

\item {} 
\sphinxAtStartPar
\sphinxcode{\sphinxupquote{\_throw\_error(msg)}} \textendash{} Raises a \sphinxcode{\sphinxupquote{ValueError}} with line number and raw text.

\end{itemize}


\bigskip\hrule\bigskip



\section{Example Usage}
\label{\detokenize{parser:example-usage}}
\begin{sphinxVerbatim}[commandchars=\\\{\}]
\PYG{n}{parser} \PYG{o}{=} \PYG{n}{NetlistParser}\PYG{p}{(}\PYG{p}{)}
\PYG{n}{components}\PYG{p}{,} \PYG{n}{analyses} \PYG{o}{=} \PYG{n}{parser}\PYG{o}{.}\PYG{n}{parse}\PYG{p}{(}\PYG{l+s+s2}{\PYGZdq{}}\PYG{l+s+s2}{my\PYGZus{}circuit.txt}\PYG{l+s+s2}{\PYGZdq{}}\PYG{p}{)}
\PYG{n}{circuit} \PYG{o}{=} \PYG{n}{Circuit}\PYG{p}{(}\PYG{n}{components}\PYG{p}{)}
\PYG{n}{simulator} \PYG{o}{=} \PYG{n}{Simulator}\PYG{p}{(}\PYG{n}{circuit}\PYG{p}{,} \PYG{n}{analyses}\PYG{p}{)}
\PYG{n}{result} \PYG{o}{=} \PYG{n}{simulator}\PYG{o}{.}\PYG{n}{execute\PYGZus{}analysis}\PYG{p}{(}\PYG{n}{sensitivity}\PYG{o}{=}\PYG{k+kc}{True}\PYG{p}{)}
\end{sphinxVerbatim}

\sphinxstepscope


\chapter{Simulation Plotting Utilities}
\label{\detokenize{plotting:simulation-plotting-utilities}}\label{\detokenize{plotting:plotting-module}}\label{\detokenize{plotting::doc}}
\sphinxAtStartPar
This module contains a collection of helper functions that take a
\sphinxcode{\sphinxupquote{SimulationResult}} object and generate industry‑standard plots:
Bode (AC), transient waveforms, DC sweep responses, sensitivity curves,
and fault‑comparison figures.  All plotting is performed with Matplotlib in a
headless mode (\sphinxtitleref{Agg}) so the utilities can be used on servers or CI pipelines.


\bigskip\hrule\bigskip

\index{built\sphinxhyphen{}in function@\spxentry{built\sphinxhyphen{}in function}!make\_bode\_plot()@\spxentry{make\_bode\_plot()}}\index{make\_bode\_plot()@\spxentry{make\_bode\_plot()}!built\sphinxhyphen{}in function@\spxentry{built\sphinxhyphen{}in function}}

\begin{fulllineitems}
\phantomsection\label{\detokenize{plotting:make_bode_plot}}
\pysigstartsignatures
\pysiglinewithargsret
{\sphinxbfcode{\sphinxupquote{make\_bode\_plot}}}
{\sphinxparam{\DUrole{n}{result}}\sphinxparamcomma \sphinxparam{\DUrole{n}{output\_nodes}}\sphinxparamcomma \sphinxparam{\DUrole{n}{folder}\DUrole{o}{=}\DUrole{default_value}{\textquotesingle{}./figures/ac\textquotesingle{}}}\sphinxparamcomma \sphinxparam{\DUrole{n}{name}\DUrole{o}{=}\DUrole{default_value}{\textquotesingle{}bodeplot\textquotesingle{}}}}
{}
\pysigstopsignatures
\sphinxAtStartPar
Generates a two‑panel Bode plot (magnitude \& phase) for the specified
\sphinxtitleref{output\_nodes}.  The result’s \sphinxcode{\sphinxupquote{sweep\_axis}} is used as frequency.


\section{Parameters}
\label{\detokenize{plotting:parameters}}
\sphinxAtStartPar
result : {\hyperref[\detokenize{results:SimulationResult}]{\sphinxcrossref{\sphinxcode{\sphinxupquote{SimulationResult}}}}}
output\_nodes : list | str
\begin{quote}

\sphinxAtStartPar
Node(s) to include.
\end{quote}
\begin{description}
\sphinxlineitem{folder, name}{[}str{]}
\sphinxAtStartPar
Where to save the PNG.

\end{description}

\end{fulllineitems}



\bigskip\hrule\bigskip



\begin{fulllineitems}

\pysigstartsignatures
\pysigline
{\sphinxbfcode{\sphinxupquote{plot\_ac\_sensitivity(result,~output\_node,~target\_component,}}}
\pysigline
{\sphinxbfcode{\sphinxupquote{folder="./figures/ac",~name="ac\_sensitivity")}}}
\pysigstopsignatures
\sphinxAtStartPar
Plots the AC magnitude of \sphinxtitleref{output\_node} together with the absolute value of its
adjoint sensitivity to \sphinxtitleref{target\_component}.

\end{fulllineitems}



\bigskip\hrule\bigskip



\begin{fulllineitems}

\pysigstartsignatures
\pysigline
{\sphinxbfcode{\sphinxupquote{plot\_transient(result,~output\_nodes=None,~folder="./figures/tran",}}}
\pysigline
{\sphinxbfcode{\sphinxupquote{name="transient")}}}
\pysigstopsignatures
\sphinxAtStartPar
Produces a standard voltage‑vs‑time trace for all nodes that are not MNA
branch currents.  If \sphinxstyleemphasis{output\_nodes} is omitted the function auto‑selects all
voltage nodes.

\end{fulllineitems}



\bigskip\hrule\bigskip



\begin{fulllineitems}

\pysigstartsignatures
\pysigline
{\sphinxbfcode{\sphinxupquote{plot\_transient\_sensitivity(result,~output\_node,}}}
\pysigline
{\sphinxbfcode{\sphinxupquote{target\_component,}}}
\pysigline
{\sphinxbfcode{\sphinxupquote{folder="./figures/tran",}}}
\pysigline
{\sphinxbfcode{\sphinxupquote{name="tran\_sensitivity")}}}
\pysigstopsignatures
\sphinxAtStartPar
Shows the transient voltage of a single node together with its time‑series
sensitivity to \sphinxtitleref{target\_component}.

\end{fulllineitems}



\bigskip\hrule\bigskip



\begin{fulllineitems}

\pysigstartsignatures
\pysigline
{\sphinxbfcode{\sphinxupquote{plot\_transient\_sensitivity\_to\_radiation(result,~output\_node,}}}
\pysigline
{\sphinxbfcode{\sphinxupquote{target\_component,}}}
\pysigline
{\sphinxbfcode{\sphinxupquote{radiation,}}}
\pysigline
{\sphinxbfcode{\sphinxupquote{folder="./figures/tran",}}}
\pysigline
{\sphinxbfcode{\sphinxupquote{name="tran\_sensitivity")}}}
\pysigstopsignatures
\sphinxAtStartPar
Similar to the previous function but plots \sphinxstyleemphasis{radiation‑scaled} sensitivities
for a list of charge values.

\end{fulllineitems}



\bigskip\hrule\bigskip



\begin{fulllineitems}

\pysigstartsignatures
\pysigline
{\sphinxbfcode{\sphinxupquote{plot\_integrated\_time\_series(result,~output\_node,}}}
\pysigline
{\sphinxbfcode{\sphinxupquote{target\_component,}}}
\pysigline
{\sphinxbfcode{\sphinxupquote{folder="./figures/tran",}}}
\pysigline
{\sphinxbfcode{\sphinxupquote{name="tran\_sensitivity")}}}
\pysigstopsignatures
\sphinxAtStartPar
Same as \sphinxcode{\sphinxupquote{plot\_transient\_sensitivity()}} \textendash{} retained for API compatibility.

\end{fulllineitems}



\bigskip\hrule\bigskip



\begin{fulllineitems}

\pysigstartsignatures
\pysigline
{\sphinxbfcode{\sphinxupquote{plot\_dc\_sweep(result,~output\_nodes=None,}}}
\pysigline
{\sphinxbfcode{\sphinxupquote{folder="./figures/dc",~name="dc\_sweep")}}}
\pysigstopsignatures
\sphinxAtStartPar
Plots the DC sweep response (voltage or current) versus the sweep source
voltage.  Auto‑selects all non‑MNA nodes if \sphinxstyleemphasis{output\_nodes} is omitted.

\end{fulllineitems}



\bigskip\hrule\bigskip



\begin{fulllineitems}

\pysigstartsignatures
\pysigline
{\sphinxbfcode{\sphinxupquote{plot\_dc\_sensitivity(result,~output\_node,}}}
\pysigline
{\sphinxbfcode{\sphinxupquote{target\_component,}}}
\pysigline
{\sphinxbfcode{\sphinxupquote{folder="./figures/dc",}}}
\pysigline
{\sphinxbfcode{\sphinxupquote{name="dc\_sensitivity")}}}
\pysigstopsignatures
\sphinxAtStartPar
Plots the DC sweep curve of a node together with its sensitivity series to
\sphinxtitleref{target\_component}.

\end{fulllineitems}



\bigskip\hrule\bigskip

\index{built\sphinxhyphen{}in function@\spxentry{built\sphinxhyphen{}in function}!print\_solution()@\spxentry{print\_solution()}}\index{print\_solution()@\spxentry{print\_solution()}!built\sphinxhyphen{}in function@\spxentry{built\sphinxhyphen{}in function}}

\begin{fulllineitems}
\phantomsection\label{\detokenize{plotting:print_solution}}
\pysigstartsignatures
\pysiglinewithargsret
{\sphinxbfcode{\sphinxupquote{print\_solution}}}
{\sphinxparam{\DUrole{n}{result}}}
{}
\pysigstopsignatures
\sphinxAtStartPar
Prints a neatly formatted table of node voltages and branch currents.
Handles both DC/OP (real) and AC (complex magnitude \& phase).

\end{fulllineitems}



\bigskip\hrule\bigskip



\begin{fulllineitems}

\pysigstartsignatures
\pysigline
{\sphinxbfcode{\sphinxupquote{plot\_fault\_comparison(circuit,~result,~output\_node,}}}
\pysigline
{\sphinxbfcode{\sphinxupquote{threshold\_results,~delta\_F,}}}
\pysigline
{\sphinxbfcode{\sphinxupquote{folder="./figures/fault",}}}
\pysigline
{\sphinxbfcode{\sphinxupquote{name="fault\_comparison")}}}
\pysigstopsignatures
\sphinxAtStartPar
Creates a fault‑comparison plot showing the nominal response, tolerance band,
and the first short/open fault responses.  Uses \sphinxtitleref{build\_xi} and
\sphinxtitleref{compute\_large\_change} from \sphinxcode{\sphinxupquote{applications.large\_change\_sensitivity}}.

\end{fulllineitems}



\bigskip\hrule\bigskip



\begin{fulllineitems}

\pysigstartsignatures
\pysigline
{\sphinxbfcode{\sphinxupquote{plot\_all\_faults(circuit,~result,~output\_node,}}}
\pysigline
{\sphinxbfcode{\sphinxupquote{threshold\_results,~delta\_F,}}}
\pysigline
{\sphinxbfcode{\sphinxupquote{folder="./figures/fault",}}}
\pysigline
{\sphinxbfcode{\sphinxupquote{name="all\_faults")}}}
\pysigstopsignatures
\sphinxAtStartPar
Extends the previous function by plotting \sphinxstyleemphasis{every} short and open fault
response found in \sphinxtitleref{threshold\_results}.  Each fault is shown as a separate curve.

\end{fulllineitems}



\bigskip\hrule\bigskip


\sphinxAtStartPar
Example usage

\begin{sphinxVerbatim}[commandchars=\\\{\}]
\PYG{k+kn}{from}\PYG{+w}{ }\PYG{n+nn}{plotting}\PYG{+w}{ }\PYG{k+kn}{import} \PYG{n}{make\PYGZus{}bode\PYGZus{}plot}\PYG{p}{,} \PYG{n}{plot\PYGZus{}transient}
\PYG{n}{result} \PYG{o}{=} \PYG{n}{run\PYGZus{}simulation\PYGZus{}core}\PYG{p}{(}\PYG{l+s+s2}{\PYGZdq{}}\PYG{l+s+s2}{my\PYGZus{}netlist.txt}\PYG{l+s+s2}{\PYGZdq{}}\PYG{p}{)}
\PYG{n}{make\PYGZus{}bode\PYGZus{}plot}\PYG{p}{(}\PYG{n}{result}\PYG{p}{,} \PYG{n}{output\PYGZus{}nodes}\PYG{o}{=}\PYG{p}{[}\PYG{l+s+s2}{\PYGZdq{}}\PYG{l+s+s2}{out}\PYG{l+s+s2}{\PYGZdq{}}\PYG{p}{]}\PYG{p}{)}
\PYG{n}{plot\PYGZus{}transient}\PYG{p}{(}\PYG{n}{result}\PYG{p}{)}
\end{sphinxVerbatim}

\sphinxstepscope


\chapter{Fault Analysis Utilities}
\label{\detokenize{fault_analysis:fault-analysis-utilities}}\label{\detokenize{fault_analysis:fault-analysis-module}}\label{\detokenize{fault_analysis::doc}}
\sphinxAtStartPar
This module implements functions to rank component sensitivities, generate fault
tables (shorts \& opens), compute resistance thresholds using large‑change
sensitivity formulas, and pretty‑print the results.  The routines operate on a
\sphinxcode{\sphinxupquote{SimulationResult}} that already contains an
\sphinxcode{\sphinxupquote{SensitivityData}} tensor.


\bigskip\hrule\bigskip

\index{built\sphinxhyphen{}in function@\spxentry{built\sphinxhyphen{}in function}!rank\_component\_sensitivities()@\spxentry{rank\_component\_sensitivities()}}\index{rank\_component\_sensitivities()@\spxentry{rank\_component\_sensitivities()}!built\sphinxhyphen{}in function@\spxentry{built\sphinxhyphen{}in function}}

\begin{fulllineitems}
\phantomsection\label{\detokenize{fault_analysis:rank_component_sensitivities}}
\pysigstartsignatures
\pysiglinewithargsret
{\sphinxbfcode{\sphinxupquote{rank\_component\_sensitivities}}}
{\sphinxparam{\DUrole{n}{result}}}
{}
\pysigstopsignatures
\sphinxAtStartPar
Rank order component sensitivities from largest to smallest.


\section{Parameters}
\label{\detokenize{fault_analysis:parameters}}
\sphinxAtStartPar
result : {\hyperref[\detokenize{results:SimulationResult}]{\sphinxcrossref{\sphinxcode{\sphinxupquote{SimulationResult}}}}}


\section{Returns}
\label{\detokenize{fault_analysis:returns}}
\sphinxAtStartPar
list{[}dict{]} | None

\sphinxAtStartPar
Prints a table of the top ten sensitivities and returns the full sorted list.

\end{fulllineitems}



\bigskip\hrule\bigskip

\index{built\sphinxhyphen{}in function@\spxentry{built\sphinxhyphen{}in function}!perform\_global\_ranking()@\spxentry{perform\_global\_ranking()}}\index{perform\_global\_ranking()@\spxentry{perform\_global\_ranking()}!built\sphinxhyphen{}in function@\spxentry{built\sphinxhyphen{}in function}}

\begin{fulllineitems}
\phantomsection\label{\detokenize{fault_analysis:perform_global_ranking}}
\pysigstartsignatures
\pysiglinewithargsret
{\sphinxbfcode{\sphinxupquote{perform\_global\_ranking}}}
{\sphinxparam{\DUrole{n}{circuit}}\sphinxparamcomma \sphinxparam{\DUrole{n}{result}}}
{}
\pysigstopsignatures
\sphinxAtStartPar
Generate a fault table ranking both open and short faults for each output node
using adjoint sensitivity data.


\section{Parameters}
\label{\detokenize{fault_analysis:id1}}
\sphinxAtStartPar
circuit : {\hyperref[\detokenize{circuit:Circuit}]{\sphinxcrossref{\sphinxcode{\sphinxupquote{Circuit}}}}}
result : {\hyperref[\detokenize{results:SimulationResult}]{\sphinxcrossref{\sphinxcode{\sphinxupquote{SimulationResult}}}}}


\section{Returns}
\label{\detokenize{fault_analysis:id2}}
\sphinxAtStartPar
dict | None

\sphinxAtStartPar
The dictionary keys are output nodes; each value is a list of fault dictionaries
sorted by absolute impact.  Each fault entry contains:
\begin{itemize}
\item {} 
\sphinxAtStartPar
\sphinxcode{\sphinxupquote{type}} \textendash{} \sphinxcode{\sphinxupquote{"Open"}} or \sphinxcode{\sphinxupquote{"Short"}}

\item {} 
\sphinxAtStartPar
\sphinxcode{\sphinxupquote{location}} \textendash{} component name or node pair

\item {} 
\sphinxAtStartPar
\sphinxcode{\sphinxupquote{impact}} \textendash{} peak magnitude of the sensitivity waveform

\item {} 
\sphinxAtStartPar
\sphinxcode{\sphinxupquote{time}} \textendash{} sweep step at which the impact occurs

\end{itemize}

\end{fulllineitems}



\bigskip\hrule\bigskip



\begin{fulllineitems}

\pysigstartsignatures
\pysigline
{\sphinxbfcode{\sphinxupquote{compute\_fault\_thresholds(circuit,~result,~output\_node,}}}
\pysigline
{\sphinxbfcode{\sphinxupquote{delta\_F,~top\_n=10,~r\_sweep\_shorts=None,}}}
\pysigline
{\sphinxbfcode{\sphinxupquote{r\_sweep\_opens=None,~dt=None,~method=\textquotesingle{}TR\textquotesingle{})}}}
\pysigstopsignatures
\sphinxAtStartPar
Compute fault resistance thresholds using the large‑change sensitivity
formula.  For each of the top faults returned by
{\hyperref[\detokenize{fault_analysis:perform_global_ranking}]{\sphinxcrossref{\sphinxcode{\sphinxupquote{perform\_global\_ranking()}}}}}, a sweep over short/open resistances is performed.
The smallest/​largest resistance that drives the output change beyond
\sphinxcode{\sphinxupquote{delta\_F}} volts is recorded.


\section{Parameters}
\label{\detokenize{fault_analysis:id3}}\begin{description}
\sphinxlineitem{circuit, result, output\_node, delta\_F}
\sphinxAtStartPar
As described above.

\end{description}

\sphinxAtStartPar
top\_n : int, optional \textendash{} number of faults to process for each type.
r\_sweep\_shorts, r\_sweep\_opens : np.ndarray, optional
\begin{quote}

\sphinxAtStartPar
Resistance sweep arrays; defaults to a wide log‑space.
\end{quote}
\begin{description}
\sphinxlineitem{dt}{[}float, optional{]}
\sphinxAtStartPar
Time step (used for companion model conductances).

\sphinxlineitem{method}{[}str, optional{]}
\sphinxAtStartPar
Integration method (‘TR’ or ‘BE’).

\end{description}


\section{Returns}
\label{\detokenize{fault_analysis:id4}}
\sphinxAtStartPar
dict

\sphinxAtStartPar
The returned dictionary contains:
\begin{itemize}
\item {} 
\sphinxAtStartPar
\sphinxcode{\sphinxupquote{shorts}} \textendash{} list of short fault threshold dictionaries.

\item {} 
\sphinxAtStartPar
\sphinxcode{\sphinxupquote{opens}}  \textendash{} list of open fault threshold dictionaries.

\item {} 
\sphinxAtStartPar
\sphinxcode{\sphinxupquote{domain}} \textendash{} string indicating the analysis domain (time, frequency, etc.).

\end{itemize}

\end{fulllineitems}



\bigskip\hrule\bigskip

\index{built\sphinxhyphen{}in function@\spxentry{built\sphinxhyphen{}in function}!print\_fault\_table()@\spxentry{print\_fault\_table()}}\index{print\_fault\_table()@\spxentry{print\_fault\_table()}!built\sphinxhyphen{}in function@\spxentry{built\sphinxhyphen{}in function}}

\begin{fulllineitems}
\phantomsection\label{\detokenize{fault_analysis:print_fault_table}}
\pysigstartsignatures
\pysiglinewithargsret
{\sphinxbfcode{\sphinxupquote{print\_fault\_table}}}
{\sphinxparam{\DUrole{n}{all\_fault\_tables}}\sphinxparamcomma \sphinxparam{\DUrole{n}{top\_n}\DUrole{o}{=}\DUrole{default_value}{15}}}
{}
\pysigstopsignatures
\sphinxAtStartPar
Pretty‑prints a ranking table for each output node.  Useful after
{\hyperref[\detokenize{fault_analysis:perform_global_ranking}]{\sphinxcrossref{\sphinxcode{\sphinxupquote{perform\_global\_ranking()}}}}}.

\end{fulllineitems}



\bigskip\hrule\bigskip

\index{built\sphinxhyphen{}in function@\spxentry{built\sphinxhyphen{}in function}!print\_threshold\_table()@\spxentry{print\_threshold\_table()}}\index{print\_threshold\_table()@\spxentry{print\_threshold\_table()}!built\sphinxhyphen{}in function@\spxentry{built\sphinxhyphen{}in function}}

\begin{fulllineitems}
\phantomsection\label{\detokenize{fault_analysis:print_threshold_table}}
\pysigstartsignatures
\pysiglinewithargsret
{\sphinxbfcode{\sphinxupquote{print\_threshold\_table}}}
{\sphinxparam{\DUrole{n}{threshold\_results}}}
{}
\pysigstopsignatures
\sphinxAtStartPar
Pretty‑prints the threshold tables generated by
\sphinxcode{\sphinxupquote{compute\_fault\_thresholds()}}.  It shows the worst short/open,
voltage values, ΔV, and observation windows.

\end{fulllineitems}



\bigskip\hrule\bigskip


\sphinxAtStartPar
Example usage

\begin{sphinxVerbatim}[commandchars=\\\{\}]
\PYG{k+kn}{from}\PYG{+w}{ }\PYG{n+nn}{fault\PYGZus{}analysis}\PYG{+w}{ }\PYG{k+kn}{import} \PYG{n}{perform\PYGZus{}global\PYGZus{}ranking}\PYG{p}{,} \PYG{n}{compute\PYGZus{}fault\PYGZus{}thresholds}
\PYG{n}{all\PYGZus{}faults} \PYG{o}{=} \PYG{n}{perform\PYGZus{}global\PYGZus{}ranking}\PYG{p}{(}\PYG{n}{circuit}\PYG{p}{,} \PYG{n}{result}\PYG{p}{)}
\PYG{n}{thresholds} \PYG{o}{=} \PYG{n}{compute\PYGZus{}fault\PYGZus{}thresholds}\PYG{p}{(}\PYG{n}{circuit}\PYG{p}{,} \PYG{n}{result}\PYG{p}{,} \PYG{l+s+s2}{\PYGZdq{}}\PYG{l+s+s2}{out}\PYG{l+s+s2}{\PYGZdq{}}\PYG{p}{,} \PYG{n}{delta\PYGZus{}F}\PYG{o}{=}\PYG{l+m+mf}{0.1}\PYG{p}{)}
\PYG{n}{print\PYGZus{}fault\PYGZus{}table}\PYG{p}{(}\PYG{n}{all\PYGZus{}faults}\PYG{p}{)}
\PYG{n}{print\PYGZus{}threshold\PYGZus{}table}\PYG{p}{(}\PYG{n}{thresholds}\PYG{p}{)}
\end{sphinxVerbatim}

\sphinxstepscope


\chapter{Large‑Change Sensitivity Utilities}
\label{\detokenize{large_change_sensitivity:largechange-sensitivity-utilities}}\label{\detokenize{large_change_sensitivity:large-change-sensitivity-module}}\label{\detokenize{large_change_sensitivity::doc}}
\sphinxAtStartPar
This module implements the Sherman‑Morrison / Kron’s formula for evaluating the exact effect of adding a resistance between two nodes in a circuit.
It requires only the LU factorisation of the Y matrix and the nominal solution vector.


\bigskip\hrule\bigskip

\index{built\sphinxhyphen{}in function@\spxentry{built\sphinxhyphen{}in function}!build\_xi()@\spxentry{build\_xi()}}\index{build\_xi()@\spxentry{build\_xi()}!built\sphinxhyphen{}in function@\spxentry{built\sphinxhyphen{}in function}}

\begin{fulllineitems}
\phantomsection\label{\detokenize{large_change_sensitivity:build_xi}}
\pysigstartsignatures
\pysiglinewithargsret
{\sphinxbfcode{\sphinxupquote{build\_xi}}}
{\sphinxparam{\DUrole{n}{n}}\sphinxparamcomma \sphinxparam{\DUrole{n}{idx\_k}}\sphinxparamcomma \sphinxparam{\DUrole{n}{idx\_l}}}
{}
\pysigstopsignatures
\sphinxAtStartPar
Build the connection vector \sphinxstylestrong{xi\_kl}.


\section{Parameters}
\label{\detokenize{large_change_sensitivity:parameters}}\begin{description}
\sphinxlineitem{n}{[}int{]}
\sphinxAtStartPar
Dimension of the Y matrix (number of nodes/branches).

\sphinxlineitem{idx\_k, idx\_l}{[}int | None{]}
\sphinxAtStartPar
Indices of the two nodes between which the fault is added.
\sphinxcode{\sphinxupquote{None}} may be passed for ground or virtual nodes.

\end{description}


\section{Returns}
\label{\detokenize{large_change_sensitivity:returns}}\begin{description}
\sphinxlineitem{np.ndarray}
\sphinxAtStartPar
Vector with +1 at \sphinxtitleref{idx\_k}, \sphinxhyphen{}1 at \sphinxtitleref{idx\_l}, zeros elsewhere.

\end{description}

\end{fulllineitems}



\bigskip\hrule\bigskip

\index{built\sphinxhyphen{}in function@\spxentry{built\sphinxhyphen{}in function}!compute\_large\_change()@\spxentry{compute\_large\_change()}}\index{compute\_large\_change()@\spxentry{compute\_large\_change()}!built\sphinxhyphen{}in function@\spxentry{built\sphinxhyphen{}in function}}

\begin{fulllineitems}
\phantomsection\label{\detokenize{large_change_sensitivity:compute_large_change}}
\pysigstartsignatures
\pysiglinewithargsret
{\sphinxbfcode{\sphinxupquote{compute\_large\_change}}}
{\sphinxparam{\DUrole{n}{lu}}\sphinxparamcomma \sphinxparam{\DUrole{n}{xi\_kl}}\sphinxparamcomma \sphinxparam{\DUrole{n}{v}}\sphinxparamcomma \sphinxparam{\DUrole{n}{R}}\sphinxparamcomma \sphinxparam{\DUrole{n}{output\_idx}}}
{}
\pysigstopsignatures
\sphinxAtStartPar
Compute the exact change in an output voltage when a resistance \sphinxstylestrong{R}
is inserted between two nodes.

\sphinxAtStartPar
The formula used:
\begin{equation*}
\begin{split}\Delta V_{\text{out}}
  = -\frac{v_{oc}}{\,R + R_{TH}\,}
    \bigl(Y^{-1} \xi_{kl}\bigr)_{\text{output}}\end{split}
\end{equation*}
\sphinxAtStartPar
where
\begin{itemize}
\item {} 
\sphinxAtStartPar
\(v_{oc} = \xi_{kl}^{T} v\) \textendash{} open‑circuit voltage between the nodes.

\item {} 
\sphinxAtStartPar
\(R_{TH} = \xi_{kl}^{T} Y^{-1} \xi_{kl}\) \textendash{} Thevenin resistance seen from
the two nodes.

\item {} 
\sphinxAtStartPar
\(Y^{-1}\xi_{kl}\) is obtained by a single forward solve using the supplied
LU factorisation.

\end{itemize}


\section{Parameters}
\label{\detokenize{large_change_sensitivity:id1}}\begin{description}
\sphinxlineitem{lu}{[}LU factorisation object{]}
\sphinxAtStartPar
Result of a prior linear solve (e.g. SciPy or custom).

\sphinxlineitem{xi\_kl}{[}np.ndarray{]}
\sphinxAtStartPar
Connection vector for the fault branch.

\sphinxlineitem{v}{[}np.ndarray{]}
\sphinxAtStartPar
Nominal solution vector at this time step.

\sphinxlineitem{R}{[}float{]}
\sphinxAtStartPar
Value of the added resistance.

\sphinxlineitem{output\_idx}{[}int{]}
\sphinxAtStartPar
Matrix index of the node whose voltage change is required.

\end{description}


\section{Returns}
\label{\detokenize{large_change_sensitivity:id2}}\begin{description}
\sphinxlineitem{float}
\sphinxAtStartPar
The exact change in the output voltage, \sphinxtitleref{ΔV\_out}.

\end{description}

\end{fulllineitems}



\bigskip\hrule\bigskip


\sphinxAtStartPar
Example usage

\begin{sphinxVerbatim}[commandchars=\\\{\}]
\PYG{k+kn}{from}\PYG{+w}{ }\PYG{n+nn}{large\PYGZus{}change\PYGZus{}sensitivity}\PYG{+w}{ }\PYG{k+kn}{import} \PYG{n}{build\PYGZus{}xi}\PYG{p}{,} \PYG{n}{compute\PYGZus{}large\PYGZus{}change}

\PYG{n}{xi} \PYG{o}{=} \PYG{n}{build\PYGZus{}xi}\PYG{p}{(}\PYG{n}{n}\PYG{o}{=}\PYG{l+m+mi}{10}\PYG{p}{,} \PYG{n}{idx\PYGZus{}k}\PYG{o}{=}\PYG{l+m+mi}{2}\PYG{p}{,} \PYG{n}{idx\PYGZus{}l}\PYG{o}{=}\PYG{l+m+mi}{5}\PYG{p}{)}
\PYG{n}{delta\PYGZus{}v} \PYG{o}{=} \PYG{n}{compute\PYGZus{}large\PYGZus{}change}\PYG{p}{(}\PYG{n}{lu}\PYG{p}{,} \PYG{n}{xi}\PYG{p}{,} \PYG{n}{v\PYGZus{}nominal}\PYG{p}{,} \PYG{n}{R}\PYG{o}{=}\PYG{l+m+mf}{100.0}\PYG{p}{,} \PYG{n}{output\PYGZus{}idx}\PYG{o}{=}\PYG{l+m+mi}{3}\PYG{p}{)}
\end{sphinxVerbatim}

\sphinxstepscope


\chapter{Circuit Component Base Class}
\label{\detokenize{component_base:circuit-component-base-class}}\label{\detokenize{component_base:component-base-module}}\label{\detokenize{component_base::doc}}
\sphinxAtStartPar
This module defines the abstract {\hyperref[\detokenize{component_base:Component}]{\sphinxcrossref{\sphinxcode{\sphinxupquote{Component}}}}} class that all circuit elements
inherit from.  It specifies the interface for matrix stamping, sensitivity
calculations, transient state updates and adjoint operations.


\bigskip\hrule\bigskip

\index{Component (built\sphinxhyphen{}in class)@\spxentry{Component}\spxextra{built\sphinxhyphen{}in class}}

\begin{fulllineitems}
\phantomsection\label{\detokenize{component_base:Component}}
\pysigstartsignatures
\pysiglinewithargsret
{\sphinxbfcode{\sphinxupquote{\DUrole{k}{class}\DUrole{w}{ }}}\sphinxbfcode{\sphinxupquote{Component}}}
{\sphinxparam{\DUrole{n}{name}}\sphinxparamcomma \sphinxparam{\DUrole{n}{data\_dict}}}
{}
\pysigstopsignatures
\sphinxAtStartPar
Abstract base class for all SPICE components.


\section{Attributes}
\label{\detokenize{component_base:attributes}}\begin{description}
\sphinxlineitem{name}{[}str{]}
\sphinxAtStartPar
Netlist identifier (e.g. \sphinxcode{\sphinxupquote{\textquotesingle{}R1\textquotesingle{}}}).

\sphinxlineitem{type}{[}str{]}
\sphinxAtStartPar
Component type character as parsed from the netlist.

\sphinxlineitem{value}{[}float{]}
\sphinxAtStartPar
Default component value (if applicable).

\sphinxlineitem{data}{[}dict{]}
\sphinxAtStartPar
Raw dictionary of parameters produced by {\hyperref[\detokenize{parser:NetlistParser}]{\sphinxcrossref{\sphinxcode{\sphinxupquote{NetlistParser}}}}}.

\sphinxlineitem{idx\_1, idx\_2}{[}int | None{]}
\sphinxAtStartPar
Matrix indices for the two terminals.  Set by \sphinxcode{\sphinxupquote{bind\_nodes()}}.

\end{description}


\section{Class Attribute}
\label{\detokenize{component_base:class-attribute}}\begin{description}
\sphinxlineitem{IS\_NONLINEAR}{[}bool{]}
\sphinxAtStartPar
\sphinxcode{\sphinxupquote{True}} if the component requires Newton‑Raphson iteration.

\end{description}

\end{fulllineitems}



\bigskip\hrule\bigskip



\bigskip\hrule\bigskip

\index{differentiable\_params (Component property)@\spxentry{differentiable\_params}\spxextra{Component property}}

\begin{fulllineitems}
\phantomsection\label{\detokenize{component_base:Component.differentiable_params}}
\pysigstartsignatures
\pysigline
{\sphinxbfcode{\sphinxupquote{\DUrole{k}{property}\DUrole{w}{ }}}\sphinxcode{\sphinxupquote{Component.}}\sphinxbfcode{\sphinxupquote{differentiable\_params}}}
\pysigstopsignatures
\sphinxAtStartPar
Returns a list of parameter names that the component can produce sensitivities for.
The default implementation returns \sphinxcode{\sphinxupquote{{[}self.name{]}}}.

\end{fulllineitems}



\bigskip\hrule\bigskip



\section{Polymorphic Matrix‑Stamping Methods}
\label{\detokenize{component_base:polymorphic-matrixstamping-methods}}
\sphinxAtStartPar
All components must implement the following methods.  They are called by the various simulation engines.


\bigskip\hrule\bigskip



\section{Adjoint Sensitivity Methods}
\label{\detokenize{component_base:adjoint-sensitivity-methods}}

\bigskip\hrule\bigskip



\section{Helper Method}
\label{\detokenize{component_base:helper-method}}

\bigskip\hrule\bigskip


\sphinxAtStartPar
Example subclass (Resistor)

\begin{sphinxVerbatim}[commandchars=\\\{\}]
\PYG{k}{class}\PYG{+w}{ }\PYG{n+nc}{Resistor}\PYG{p}{(}\PYG{n}{Component}\PYG{p}{)}\PYG{p}{:}
    \PYG{n}{IS\PYGZus{}NONLINEAR} \PYG{o}{=} \PYG{k+kc}{False}
    \PYG{k}{def}\PYG{+w}{ }\PYG{n+nf}{stamp\PYGZus{}static}\PYG{p}{(}\PYG{n+nb+bp}{self}\PYG{p}{,} \PYG{n}{Y}\PYG{p}{)}\PYG{p}{:}
        \PYG{n}{g} \PYG{o}{=} \PYG{l+m+mf}{1.0} \PYG{o}{/} \PYG{n+nb+bp}{self}\PYG{o}{.}\PYG{n}{value}
        \PYG{n}{i}\PYG{p}{,} \PYG{n}{j} \PYG{o}{=} \PYG{n+nb+bp}{self}\PYG{o}{.}\PYG{n}{idx\PYGZus{}1}\PYG{p}{,} \PYG{n+nb+bp}{self}\PYG{o}{.}\PYG{n}{idx\PYGZus{}2}
        \PYG{k}{if} \PYG{n}{i} \PYG{o+ow}{is} \PYG{o+ow}{not} \PYG{k+kc}{None}\PYG{p}{:}
            \PYG{n}{Y}\PYG{p}{[}\PYG{n}{i}\PYG{p}{,}\PYG{n}{i}\PYG{p}{]} \PYG{o}{+}\PYG{o}{=} \PYG{n}{g}
            \PYG{n}{Y}\PYG{p}{[}\PYG{n}{j}\PYG{p}{,}\PYG{n}{j}\PYG{p}{]} \PYG{o}{+}\PYG{o}{=} \PYG{n}{g}
        \PYG{k}{if} \PYG{n}{i} \PYG{o+ow}{is} \PYG{o+ow}{not} \PYG{k+kc}{None} \PYG{o+ow}{and} \PYG{n}{j} \PYG{o+ow}{is} \PYG{o+ow}{not} \PYG{k+kc}{None}\PYG{p}{:}
            \PYG{n}{Y}\PYG{p}{[}\PYG{n}{i}\PYG{p}{,}\PYG{n}{j}\PYG{p}{]} \PYG{o}{\PYGZhy{}}\PYG{o}{=} \PYG{n}{g}
            \PYG{n}{Y}\PYG{p}{[}\PYG{n}{j}\PYG{p}{,}\PYG{n}{i}\PYG{p}{]} \PYG{o}{\PYGZhy{}}\PYG{o}{=} \PYG{n}{g}
\end{sphinxVerbatim}

\sphinxstepscope


\chapter{Capacitor Component}
\label{\detokenize{component_capacitor:capacitor-component}}\label{\detokenize{component_capacitor:id1}}\label{\detokenize{component_capacitor::doc}}
\sphinxAtStartPar
This module implements a dynamic capacitor that participates in AC, transient,
and adjoint sensitivity calculations.  It inherits from \sphinxcode{\sphinxupquote{Component}}.


\bigskip\hrule\bigskip

\index{Capacitor (built\sphinxhyphen{}in class)@\spxentry{Capacitor}\spxextra{built\sphinxhyphen{}in class}}

\begin{fulllineitems}
\phantomsection\label{\detokenize{component_capacitor:Capacitor}}
\pysigstartsignatures
\pysiglinewithargsret
{\sphinxbfcode{\sphinxupquote{\DUrole{k}{class}\DUrole{w}{ }}}\sphinxbfcode{\sphinxupquote{Capacitor}}}
{\sphinxparam{\DUrole{n}{name}}\sphinxparamcomma \sphinxparam{\DUrole{n}{data\_dict}}}
{}
\pysigstopsignatures
\sphinxAtStartPar
Dynamic capacitor component.

\sphinxAtStartPar
Inherits all attributes and abstract methods from {\hyperref[\detokenize{component_base:Component}]{\sphinxcrossref{\sphinxcode{\sphinxupquote{Component}}}}}.
The constructor simply forwards to the base class; the key behaviour
is implemented in the overridden methods below.

\end{fulllineitems}



\bigskip\hrule\bigskip



\section{Overridden Methods}
\label{\detokenize{component_capacitor:overridden-methods}}

\bigskip\hrule\bigskip


\sphinxAtStartPar
Example Usage

\begin{sphinxVerbatim}[commandchars=\\\{\}]
\PYG{k+kn}{from}\PYG{+w}{ }\PYG{n+nn}{components}\PYG{n+nn}{.}\PYG{n+nn}{capacitor}\PYG{+w}{ }\PYG{k+kn}{import} \PYG{n}{Capacitor}

\PYG{n}{cap} \PYG{o}{=} \PYG{n}{Capacitor}\PYG{p}{(}\PYG{l+s+s1}{\PYGZsq{}}\PYG{l+s+s1}{C1}\PYG{l+s+s1}{\PYGZsq{}}\PYG{p}{,} \PYG{p}{\PYGZob{}}\PYG{l+s+s1}{\PYGZsq{}}\PYG{l+s+s1}{type}\PYG{l+s+s1}{\PYGZsq{}}\PYG{p}{:} \PYG{l+s+s1}{\PYGZsq{}}\PYG{l+s+s1}{C}\PYG{l+s+s1}{\PYGZsq{}}\PYG{p}{,} \PYG{l+s+s1}{\PYGZsq{}}\PYG{l+s+s1}{n1}\PYG{l+s+s1}{\PYGZsq{}}\PYG{p}{:} \PYG{l+m+mi}{1}\PYG{p}{,} \PYG{l+s+s1}{\PYGZsq{}}\PYG{l+s+s1}{n2}\PYG{l+s+s1}{\PYGZsq{}}\PYG{p}{:} \PYG{l+m+mi}{0}\PYG{p}{,} \PYG{l+s+s1}{\PYGZsq{}}\PYG{l+s+s1}{value}\PYG{l+s+s1}{\PYGZsq{}}\PYG{p}{:} \PYG{l+m+mf}{1e\PYGZhy{}6}\PYG{p}{\PYGZcb{}}\PYG{p}{)}
\PYG{n}{cap}\PYG{o}{.}\PYG{n}{bind\PYGZus{}nodes}\PYG{p}{(}\PYG{n}{node\PYGZus{}map}\PYG{p}{)}
\PYG{n}{cap}\PYG{o}{.}\PYG{n}{stamp\PYGZus{}ac}\PYG{p}{(}\PYG{n}{Y}\PYG{p}{,} \PYG{n}{sources}\PYG{p}{,} \PYG{n}{w}\PYG{o}{=}\PYG{l+m+mi}{2}\PYG{o}{*}\PYG{n}{np}\PYG{o}{.}\PYG{n}{pi}\PYG{o}{*}\PYG{l+m+mi}{1000}\PYG{p}{)}
\end{sphinxVerbatim}


\chapter{Indices and tables}
\label{\detokenize{index:indices-and-tables}}\begin{itemize}
\item {} 
\sphinxAtStartPar
\DUrole{xref}{\DUrole{std}{\DUrole{std-ref}{genindex}}}

\item {} 
\sphinxAtStartPar
\DUrole{xref}{\DUrole{std}{\DUrole{std-ref}{modindex}}}

\item {} 
\sphinxAtStartPar
\DUrole{xref}{\DUrole{std}{\DUrole{std-ref}{search}}}

\end{itemize}



\renewcommand{\indexname}{Index}
\printindex
\end{document}